% ============================================================
%  第五卷 第十二章 涨落（前半）（§110–§118 + 习题）
%  底本：高教社中文版第5版（束仁贵/束莼 译，郑伟谋 校）
%  OCR 底稿 full_merged.md 行 11515–12778；对照 原书转换/pages/p-308…p-334.png 核验
%  插图：本区间（§110–§118）原书无插图。
%  正文圈码 26 处原样保留（§110 四处、§111 三处、§112 四处、§114 两处、
%  §115 一处、§116 四处、§117 五处、§118 三处），脚注文本由脚注代理后续转录。
%
%  已修正的 OCR/原书错误（除注明外均对照原书页面 PNG 核验）：
%   [§110] (110.1) 按原书还原为 w(x)=常数·e^{S(x)}（OCR 乱码，p293）；
%          (110.5) w⟨x⟩→w(x)；(110.7) 补下标 x=0（p295）；\mathbb{O}→①；
%          "显然，x 这个量…"（OCR "显然$,x$"）。
%   [§111] (111.4) δ_lm（OCR 8_lm）；"两个行列式 a=|a_ik|"（OCR 粗体 * 乱码）；
%          ⟨X_iX_k⟩=β_il⟨x_lX_k⟩=β_ilδ_lk（OCR β_ι）；
%          "热力学共轭量"后 OCR 乱码 \mathfrak{T} 实为脚注圈码①（p296）；
%          "(111.7a)"系脚注①内公式编号窜入正文（正文积公式无编号），按 p296 删；
%          w_12 之"常数·"按 p297 还原，指数末项符号 OCR 作"="正为"−"；
%          逆矩阵式 β'^{-1}_12=β'_12/(β'^2_12−β_11β_22) 系脚注①内容窜入正文，
%          按 p297 从正文删除（脚注文本归脚注代理）。
%   [★§112] (112.3) 原书排印误作 ΔPΔG（p299）——按前式 ½(ΔSΔT−ΔPΔV) 之推导
%          正为 ΔPΔV；
%   [★§112] 习题3 原书排印误作 (∂P/∂V)_V（p303）——按 T、V 涨落统计独立，
%          ⟨ΔTΔP⟩=(∂P/∂T)_V·⟨(ΔT)²⟩ 正为 (∂P/∂T)_V；
%          (112.5)(112.9) 及习题4–6 补平均角标 ⟨⟩（OCR 脱）；习题4 原书即 −T
%          （p303），OCR 脱负号，已补；习题7 R_min=½mg·lφ²（OCR 乱码）；
%          习题8 (l−x)²（OCR "l−x²"）；"在恒定的 N 下取的"（OCR "恒定T的N"，p301）；
%          S_t、ΔS_t、ΔF_t 之下标 t（OCR s_i、s,、S_∥、S_r、F_1、F_ι、F_∥）。
%   [§113] (113.2) ∂n̄_k/∂μ（OCR 脱平均线）；"相邻各 n̄_k"（OCR n̄_★）；
%          T_0（OCR T_\mathfrak{d}）；"代入"（OCR "代人"）。
%   [§114] "我们可以令 N_0!/(N_0−N)!≈N_0^N"（OCR 碎裂，按 p307 还原）；
%          上标 N、N_0（OCR ν、ν_0）；Σε_k（OCR Σσ_k）；e^{Ω/T}（OCR e^{πτ}）；
%          "常数·N̄^N/N!"与"常数=e^{−N̄}①"按 p308 还原（OCR 乱码）；
%          (114.1) 按 p306 照排"·dN"之点乘。
%   [§115] R_min=ΔΦ−μ′Δn（OCR φ→Φ）；下标 P,T（OCR "P. T"）；
%          "要计算均方值 ⟨(Δc²)⟩"照原书排（p309；角标位置与(115.3)不一致，原书如此）。
%   [★§116] (116.2) 原书排印误作 N̄δ(r_2−r_1)（p310）——按其后"带δ函数的项给出
%          一阶量 n̄ΔV"及 ∬δdV_1dV_2=V 自洽，正为 n̄δ(r_2−r_1)；
%          (116.1) 右端第一项按原书排长横线 \overline{n_1n_2}（p309）；
%          δ 函数（OCR "S"）；(116.7)(116.8) 之 ΔF_t（同上）；
%          习题 ν(r)=n̄[e^{−U(r)/T}−1]、E=E_id+(N²/2V)∫Ue^{−U/T}dV
%          （OCR e^{−U/τ}、e^{−υτ}）。
%   [§117] 算符记号统一为 â⁺、ψ*（OCR a^*/a^†/a^+、ψ^+/ψ^* 混排，按 p312–p313）；
%          (117.2) e^{ip·r/ħ}（OCR /h）；(117.8) 原书脱右端 |·|² 之平方（p314），
%          按(117.5)(117.7)及(117.11)反推补上；
%          (117.11) 后式分母 p_F r⁴（OCR r²）、cos²(p_F r/ħ)（OCR 脱 r）；
%          习题1 n̄_{p+ħk}（OCR n_{p+nk}）、"中心移了 ħk"（OCR h）；
%          习题2 x=p_F(p−p_F)/(mT)、λ=mT/(ħp_F)、ε_F（OCR ε_r）；
%          习题3 测度 (2πħ)³——原书误排 (2πħ²)³（p316），按其结果式
%          I=mT/(2πħ²r) 反推改正；
%   [★§117] 习题3"求得①"之脚注（傅里叶变换公式两则，κ 积分）窜入正文
%          （OCR 断在习题4题干之后），按 p316 从正文删除（脚注文本归脚注代理）；
%          习题4 I=e^{ip·r/ħ}n̄_p…（OCR /κ、n̄_ρ）、第二项 4π²（OCR 脱平方）、
%          ∫νdV（OCR ∫νdν）、T_0（OCR T_\mathfrak{g}、T_\mathfrak{o}）。
%   [§118] (118.4) 删 OCR 衍余之再一个 ⟩；(118.7) −λξ_x(t)（OCR λẋ_x）；
%          (118.8) e^{−λ|t|}（OCR 脱绝对值）；(118.9) ẋ=−λx+y（OCR 脱点）；
%          (118.10) δ(t)（OCR \tilde{⊗}）；"虽然 y 这个量"（OCR γ）；
%          "（即 x=0 下）"（OCR "\top"）；"给定的差值 t′−t"（OCR 乱码）；
%          x̄=0（OCR x̃=0）；"由于 x 很小我们有①"等处圈码照排。
%
%  遗留问题：
%   - 脚注仅存圈码（共 26 处，分布见上），脚注文本由脚注代理转录；
%   - §112 "它来自式(24.14) dΩ=VdP=SdT+Ndμ" 照原书排（p302）；
%     其中 "dΩ=VdP" 之符号疑有误（由 Ω=−PV 推得者应为 −VdP），且卷5已完成章节
%     之 §24 编号仅至(24.13)，(24.14) 之指称待全书装配时校核；
%   - §117 习题2 之 I=+ħ²∂/∂r[…] 与上一等号 −ħ²∂/∂r[…] 符号相左，原书如此
%     照排（p316；ν(r) 只用 |I|²，不影响最终结果）；
%   - §117 习题2 末式原书括号外衍一上标 2（sinh^{-2}(…)²），按推导应为
%     sinh^{-2}(πmTr/ħp_F)，已按物理修正；
%   - §112 习题2 题干 ⟨(ΔW)²⟩ 之 W：OCR 作 \mathbb{W}，按焓记号正为 W。
%
%  脚注转录记录（脚注代理，2026-10-12）：
%   - 全章 26 处圈码已转录为 \footnote，实际页位（PNG 页码）：§110 p308①②、
%     p309①、p310①；§111 p310②、p311①、p312①；§112 p313①、p315①②、
%     p317①；§114 p322①、p323①；§115 p323②；§116 p325①、p326①②③；
%     §117 p327①、p328①、p329①②、p331①；§118 p333①②③。
%   - 窜入正文之三处脚注内容已随脚注归位：§111 脚注①收 (111.7a) 及 dS 推演链
%     （p311）；§111 脚注①（p312）收逆矩阵式——原备案作 β′ 撇标，页面实况
%     无撇（β₁₂⁻¹=β₁₂/(β₁₂²−β₁₁β₂₂)），照原书转录；§117 习题3 脚注①（p331）
%     收傅里叶变换 κ 积分公式两则及 Δφ−κ²φ=−4πδ(r) 推演。
%   - [疑难] §115 脚注②（p323）"x₁=ΔS₁"之 S 带下标 1，原书如此（按文意应为
%     ΔS），照原书转录，留待校核。
% ============================================================
\chapter{涨落}

\section{高斯分布}

我们已经一再强调：表征处于平衡状态的宏观物体的各物理量，实际上总是在很大的精确程度上等于它们的平均值.但是对平均值的偏离无论多么小，它们总是存在的（通常说：这些量在涨落），因而发生寻求这些偏差的概率分布的问题.

我们来考虑某一闭合系统，设 $x$ 为表征整个系统或其某一部分的某个物理量（在前一种情况下，这个量不一定要像能量那样在闭合系统中严格地保持恒定）.今后，为方便起见，假设已经从 $x$ 中减去平均值 $\bar{x}$，因此在下面处处假定 $\bar{x}=0$.

在§7中所作的讨论表明：如果把系统的熵形式地看成子系统能量精确值的函数，那么函数 $\mathrm{e}^{S}$ 给出能量的概率分布（公式(7.17)）.但是很容易看出：在讨论中根本没有用到能量特有的任何性质.因此，同样的推理可导致下面的结果：$x$ 这个量的值在 $x$ 和 $x+\mathrm{d}x$ 间隔内的概率正比于 $\mathrm{e}^{S(x)}$，式中 $S(x)$ 是在形式上作为 $x$ 精确值的函数的熵.把这个概率用 $w(x)\,\mathrm{d}x$ 来表示，我们有\footnote{该公式首先由 A.爱因斯坦（1907）用于研究涨落.}：
\begin{equation}
  w(x)=\text{常数}\cdot\mathrm{e}^{S(x)}.
\end{equation}

在对该公式进一步研究以前，先考虑一下它的适用范围.得出公式(110.1)的整个推导隐含地假定了 $x$ 这个量的经典性质\footnote{当然，这并不意味着整个系统必须是经典的. 系统的（除 $x$ 以外的）其它量可以是量子性的.}.因此我们必须求出保证量子效应可以忽略的条件.

在量子力学中大家知道，在能量和某一个量 $x$ 二者的量子不确定度之间存在着关系式
\begin{equation*}
  \Delta E\,\Delta x\sim\hbar\dot{x},
\end{equation*}
式中 $\dot{x}$ 是 $x$ 这个量的经典变化率（参看第三卷§16）.

设我们所感兴趣的 $x$ 这个量具有非平衡值，且它的变化率可以用时间 $\tau$ 来表征\footnote{时间 $\tau$ 并不一定是 $x$ 这个量达到平衡值的弛豫时间，如果它以振荡的方式趋于 $\bar{x}$，则 $\tau$ 可以小于弛豫时间. 例如，如果我们所讨论的是物体的某一小部分（其线度 $\sim a$）中压强的变化，那么 $\tau$ 与波长为 $\lambda\sim a$ 的声振动的周期同数量级，即 $\tau\sim a/c$，其中 $c$ 为声速.}；这时 $\dot{x}\sim\dfrac{x}{\tau}$，因此
\begin{equation*}
  \Delta E\,\Delta x\sim\frac{\hbar x}{\tau}.
\end{equation*}
显然，$x$ 这个量只能在其量子不确定度很小的条件下：$\Delta x\ll x$，可以说具有确定值；于是
\begin{equation*}
  \Delta E\gg\frac{\hbar}{\tau}.
\end{equation*}
由此可见，能量的不确定度必须比 $\dfrac{\hbar}{\tau}$ 大得多.因而系统的熵这时具有不确定度
\begin{equation*}
  \Delta S\gg\frac{\hbar}{\tau T}.
\end{equation*}
要使公式(110.1)具有实际意义，显然必须使熵的不确定度远小于1：$\Delta S\ll1$，由此得出
\begin{equation}
  T\gg\frac{\hbar}{\tau},\qquad \tau\gg\frac{\hbar}{T}.
\end{equation}
这就是所求的条件.当温度太低或 $x$ 这个量的变化太快（$\tau$ 太小）时，涨落就不能再按热力学处理，而纯量子的涨落将起首要的作用.

现在我们回到(110.1)式.熵 $S$ 在 $x=\bar{x}=0$ 时具有极大值.因此
\begin{equation*}
  \left.\frac{\partial S}{\partial x}\right|_{x=0}=0,\qquad
  \left.\frac{\partial^{2}S}{\partial x^{2}}\right|_{x=0}<0.
\end{equation*}
涨落中的 $x$ 这个量很小.把 $S(x)$ 展开成 $x$ 的幂级数，并只取到二次项，我们得到：
\begin{equation}
  S(x)=S(0)-\frac{\beta}{2}x^{2},
\end{equation}
式中 $\beta$ 是一正常数.代入(110.1)，我们就得到下面形式的概率分布：
\begin{equation*}
  w(x)\,\mathrm{d}x=A\exp\Bigl(-\frac{\beta}{2}x^{2}\Bigr)\mathrm{d}x.
\end{equation*}
归一化常数 $A$ 由条件
\begin{equation*}
  \int w(x)\,\mathrm{d}x=1
\end{equation*}
给出.虽然我们在这里所用的 $w(x)$ 的表达式只适用于很小的 $x$，但是被积函数随 $|x|$ 的增加而很快地减小，所以积分范围可以拓展为从 $-\infty$ 到 $+\infty$ 的所有值.积分后我们得到 $A=\sqrt{\beta/2\pi}$.

因此涨落 $x$ 的不同值的概率分布由以下公式确定：
\begin{equation}
  w(x)\,\mathrm{d}x=\sqrt{\frac{\beta}{2\pi}}\,
  \exp\Bigl(-\frac{\beta}{2}x^{2}\Bigr)\mathrm{d}x.
\end{equation}
这种类型的分布称为高斯分布；它在 $x=0$ 处具有极大值，并在其两侧对称地随着 $|x|$ 的增加而很快减小.

涨落的均方值等于
\begin{equation}
  \left\langle x^{2}\right\rangle
  =\int_{-\infty}^{\infty}x^{2}w(x)\,\mathrm{d}x=\frac{1}{\beta}.
\end{equation}
所以高斯分布可以写成
\begin{equation}
  w(x)\,\mathrm{d}x=\frac{1}{\sqrt{2\pi\langle x^{2}\rangle}}\,
  \exp\Bigl(-\frac{x^{2}}{2\langle x^{2}\rangle}\Bigr)\mathrm{d}x.
\end{equation}
正如所料，$\langle x^{2}\rangle$ 愈小，则 $w(x)$ 的极大值愈尖锐.

应该指出，知道了均方值 $\langle x^{2}\rangle$，就可以求出任何函数 $\varphi(x)$ 的均方值.由于 $x$ 很小我们有\footnote{当然，这里假定了对于 $x\sim\sqrt{\langle x^{2}\rangle}$ 的取值函数 $\varphi(x)$ 改变很小，而且在 $x=0$ 处导数 $\dfrac{\mathrm{d}\varphi}{\mathrm{d}x}$ 不为零.}：
\begin{equation}
  \left\langle(\Delta\varphi)^{2}\right\rangle
  =\left(\frac{\mathrm{d}\varphi}{\mathrm{d}x}\right)^{2}_{x=0}
  \left\langle x^{2}\right\rangle.
\end{equation}

\section{多个热力学量的高斯分布}

在上节中，我们讨论了任意一个热力学量偏离其平均值的概率，而并不考虑其它量的值，亦即认为它们取值任意\footnote{这意味着：我们在§110中所用的函数 $S(x)$ 是在给定的非平衡值 $x$ 下熵可能取的最大值.}.用类似的方式也可以确定一系列热力学量同时偏离各自平均值的概率；我们用 $x_{1},x_{2},\cdots,x_{n}$ 来代表这些偏差.

我们引入熵 $S(x_{1},\cdots,x_{n})$ 作为所考虑各量 $x_{1},x_{2},\cdots,x_{n}$ 的函数，并且用(110.1)的 $w$ 把概率分布写成
\begin{equation*}
  w\,\mathrm{d}x_{1}\mathrm{d}x_{2}\cdots\mathrm{d}x_{n}.
\end{equation*}
我们把 $S$ 按 $x_{i}$ 的幂次展开；取到二次项，差 $S-S_{0}$ 可以表示为负定二次型的形式
\begin{equation*}
  S-S_{0}=-\frac{1}{2}\sum_{i,k=1}^{n}\beta_{ik}x_{i}x_{k}.
\end{equation*}
显然 $\beta_{ik}=\beta_{ki}$.在本节中以后我们将省略求和号，而总是约定对重复出现两次的角标求和（遍及从1到 $n$ 的所有值）.于是我们记：
\begin{equation}
  S-S_{0}=-\frac{1}{2}\beta_{ik}x_{i}x_{k}.
\end{equation}
把这个式子代入(110.1)，我们就得到所求的概率分布的表达式
\begin{equation}
  w=A\exp\Bigl(-\frac{1}{2}\beta_{ik}x_{i}x_{k}\Bigr).
\end{equation}
常数 $A$ 取决于归一化条件 $\int w\,\mathrm{d}x_{1}\cdots\mathrm{d}x_{n}=1$，其中（根据和§110中同样的理由）对所有 $x_{i}$ 的积分都可以从 $-\infty$ 取到 $\infty$.积分可计算如下.我们对上述各量 $x_{i}$ 作线性变换
\begin{equation}
  x_{i}=a_{ik}x'_{k},
\end{equation}
它把二次型 $\beta_{ik}x_{i}x_{k}$ 化为平方和 $x'_{i}x'_{i}$.为了满足
\begin{equation*}
  \beta_{ik}x_{i}x_{k}=x'_{i}x'_{i}\equiv x'_{i}x'_{k}\delta_{ik},
\end{equation*}
变换系数 $a_{ik}$ 应满足
\begin{equation}
  \beta_{ik}a_{il}a_{km}=\delta_{lm}.
\end{equation}
这个等式的左边量的矩阵行列式等于行列式 $\beta=|\beta_{ik}|$ 和两个行列式 $a=|a_{ik}|$ 的乘积.然而，行列式 $|\delta_{ik}|=1$，所以由上述关系式得出
\begin{equation}
  \beta a^{2}=1.
\end{equation}
从变量 $x_{i}$ 到变量 $x'_{i}$ 的线性变换的雅可比式是常量即行列式 $a$.所以变换以后归一化的积分分解成 $n$ 个相同积分的乘积，考虑到(111.5)就得到
\begin{equation*}
  Aa\left[\int_{-\infty}^{\infty}
  \exp\Bigl(-\frac{1}{2}x^{\prime 2}\Bigr)\mathrm{d}x'\right]^{n}
  =\frac{A}{\sqrt{\beta}}\,(2\pi)^{n/2}=1.
\end{equation*}
因此我们最后求出多变量高斯分布的形式为
\begin{equation}
  w=\frac{\sqrt{\beta}}{(2\pi)^{n/2}}
  \exp\Bigl(-\frac{1}{2}\beta_{ik}x_{i}x_{k}\Bigr).
\end{equation}
我们引入
\begin{equation}
  X_{i}=-\frac{\partial S}{\partial x_{i}}=\beta_{ik}x_{k},
\end{equation}
称之为量 $x_{i}$ 的热力学共轭量\footnote{值得注意，在线性关系(111.7)下，共轭性是双向的：如果同一个熵 $S$ 用 $X_i$ 各量表示，则
\begin{equation*}
  x_{i}=-\frac{\partial S}{\partial X_{i}}.\tag{111.7a}
\end{equation*}
实际上，利用(111.7)我们有
\begin{equation*}
  \mathrm{d}S=-X_{k}\,\mathrm{d}x_{k}
  =-\beta_{ki}x_{i}\,\mathrm{d}x_{k}
  =-x_{i}\,\mathrm{d}(\beta_{ik}x_{k})
  =-x_{i}\,\mathrm{d}X_{i}.
\end{equation*}}.乘积 $x_{i}X_{k}$ 的平均值由下式确定：
\begin{equation*}
  \left\langle x_{i}X_{k}\right\rangle=\frac{\sqrt{\beta}}{(2\pi)^{n/2}}
  \int\cdots\int x_{i}\beta_{kl}x_{l}
  \exp\Bigl(-\frac{1}{2}\beta_{ik}x_{i}x_{k}\Bigr)
  \mathrm{d}x_{1}\cdots\mathrm{d}x_{n}.
\end{equation*}
为了求出这个积分，我们暂时假设：平均值 $\bar{x}_{i}$ 不等于0，而等于某一有限值 $x_{i0}$.于是在(111.6)中我们必须把 $x_{i}$ 写成 $x_{i}-x_{i0}$，并且按平均值的定义我们得到：
\begin{equation*}
  \bar{x}_{i}=\frac{\sqrt{\beta}}{(2\pi)^{n/2}}
  \int\cdots\int x_{i}\exp\Bigl[-\frac{1}{2}\beta_{ik}
  (x_{i}-x_{i0})(x_{k}-x_{k0})\Bigr]\mathrm{d}x_{1}\cdots\mathrm{d}x_{n}
  =x_{i0}.
\end{equation*}
这个等式对 $x_{k0}$ 微分，随后重新令所有的 $x_{10},\cdots,x_{n0}$ 都等于0，于是在右边得到 $\delta_{ik}$，而在左边正好得所求的积分.

因此我们得到：
\begin{equation}
  \left\langle x_{i}X_{k}\right\rangle=\delta_{ik}.
\end{equation}
把(111.7)式代入上式得到：
\begin{equation*}
  \beta_{kl}\left\langle x_{l}x_{i}\right\rangle=\delta_{ik},
\end{equation*}
由此得出
\begin{equation}
  \left\langle x_{i}x_{k}\right\rangle=\beta_{ik}^{-1},
\end{equation}
式中 $\beta_{ik}^{-1}$ 是 $\beta_{ik}$ 的逆矩阵的矩阵元.

最后，我们来计算 $\langle X_{i}X_{k}\rangle$.根据(111.7)和(111.8)，我们得到
$\left\langle X_{i}X_{k}\right\rangle=\beta_{il}\left\langle x_{l}X_{k}\right\rangle
=\beta_{il}\delta_{lk}$，即
\begin{equation}
  \left\langle X_{i}X_{k}\right\rangle=\beta_{ik}.
\end{equation}

任意函数 $\varphi(x_{1},\cdots,x_{n})$ 的均方涨落也可以很容易求出.因为相对于平均值的偏差很小，所以
$\Delta\varphi=\dfrac{\partial\varphi}{\partial x_{i}}\Delta x_{i}$.在这里
$\dfrac{\partial\varphi}{\partial x_{i}}$ 记 $x_{1}=x_{2}=\cdots=0$ 处的微商值.由此得出
\begin{equation}
  \left\langle(\Delta\varphi)^{2}\right\rangle
  =\frac{\partial\varphi}{\partial x_{i}}
  \frac{\partial\varphi}{\partial x_{k}}
  \left\langle x_{i}x_{k}\right\rangle
  =\frac{\partial\varphi}{\partial x_{i}}
  \frac{\partial\varphi}{\partial x_{k}}\beta_{ik}^{-1}.
\end{equation}

如果有任意两个 $x_{i}$（譬如说是 $x_{1}$ 和 $x_{2}$）的涨落是统计独立的，那么平均值
$\langle x_{1}x_{2}\rangle$ 就等于平均值 $\bar{x}_{1}$ 和 $\bar{x}_{2}$ 之积，而因为
$\bar{x}_{1}$ 和 $\bar{x}_{2}$ 都等于0，所以 $\langle x_{1}x_{2}\rangle$ 也等于0.根据(111.9)，这表明
$\beta_{12}^{-1}=0$.不难看出，在高斯概率分布的情形下，逆定理也成立：如果
$\langle x_{1}x_{2}\rangle=0$（即 $\beta_{12}^{-1}=0$），则 $x_{1}$ 和 $x_{2}$
这两个量的涨落是统计独立的.

事实上，两个量 $x_{1}$ 和 $x_{2}$ 的概率分布 $w_{12}$，可通过把分布函数(111.6)对所有其余的
$x_{i}$ 积分而得到；这样得到的表达式形式为
\begin{equation*}
  w_{12}=\text{常数}\cdot\exp\Bigl\{-\frac{1}{2}\beta'_{11}x_{1}^{2}
  -\beta'_{12}x_{1}x_{2}-\frac{1}{2}\beta'_{22}x_{2}^{2}\Bigr\}
\end{equation*}
（式中系数 $\beta'_{ik}$ 一般不等于相应的矩阵元 $\beta_{ik}$）.把(111.9)式应用于这个分布，我们就求出：
$\left\langle x_{1}x_{2}\right\rangle=\beta_{12}^{\prime-1}$.如果
$\langle x_{1}x_{2}\rangle=0$，则 $\beta_{12}^{\prime-1}=0$.但是对于一个二秩矩阵来讲，逆矩阵元
$\beta_{12}^{\prime-1}$ 为零意味着原矩阵元 $\beta'_{12}$ 为零\footnote{对于一个二秩矩阵，我们有：
\begin{equation*}
  \beta_{12}^{-1}=\frac{\beta_{12}}{\beta_{12}^{2}-\beta_{11}\beta_{22}}.
\end{equation*}}.结果，$w_{12}$ 分解成关于
$x_{1}$ 和 $x_{2}$ 的两个独立高斯分布之积，这也就是说它们是统计独立的.

\begin{xiti}

求平均值 $\langle\exp(\alpha_{i}x_{i})\rangle$，式中 $\alpha_{i}$ 是常数，而 $x_{i}$
是服从高斯分布(111.2)的涨落量.

解：这里须计算积分
\begin{equation*}
  \left\langle\exp(\alpha_{i}x_{i})\right\rangle=A\int
  \exp\Bigl(\alpha_{i}x_{i}-\frac{1}{2}\beta_{ik}x_{i}x_{k}\Bigr)
  \mathrm{d}x_{1}\cdots\mathrm{d}x_{n}.
\end{equation*}
用变换(111.3)把被积函数的指数化为
\begin{equation*}
  \alpha_{i}x_{i}-\frac{1}{2}\beta_{ik}x_{i}x_{k}
  =\alpha_{i}a_{ik}x'_{k}-\frac{1}{2}x_{k}^{\prime 2}
  =-\frac{1}{2}(x'_{k}-\alpha_{i}a_{ik})^{2}
  +\frac{1}{2}\alpha_{i}a_{ik}\alpha_{l}a_{lk},
\end{equation*}
此后，积分给出
\begin{equation*}
  \left\langle\exp(\alpha_{i}x_{i})\right\rangle
  =\exp\Bigl(\frac{1}{2}\alpha_{i}\alpha_{l}a_{ik}a_{lk}\Bigr).
\end{equation*}
根据(111.4)有 $a_{ik}=a^{-1}_{km}\beta^{-1}_{mi}$，所以
$a_{ik}a_{lk}=\beta^{-1}_{li}$.因此，考虑到(111.9)，最后有
\begin{equation*}
  \left\langle\exp(\alpha_{i}x_{i})\right\rangle
  =\exp\Bigl(\frac{1}{2}\alpha_{i}\alpha_{k}
  \left\langle x_{i}x_{k}\right\rangle\Bigr).
\end{equation*}

\end{xiti}

\section{基本热力学量的涨落}

我们现在对物体中划出的任意一小部分计算基本热力学量的均方涨落.自然，这一小部分仍旧包含足够多的粒子.但是，在很低的温度下，这个条件比保证不出现量子涨落的条件(110.2)更弱；在这种情况下，物体这一部分的极小允许尺度就由后一条件决定\footnote{例如，对于压强的涨落来讲，条件 $\tau\gg\hbar/T$ 取 $\tau\sim a/c$ 时（参看§110第三个脚注）给出
\begin{equation*}
  a\gg\hbar c/T.
\end{equation*}}.为了避免误解起见，必须强调指出：关于量子涨落如何重要的问题，与量子效应对物质的热力学量（物态方程）有何影响的问题没有任何关系；涨落可能是经典的，但同时物态方程却可以取决于量子力学公式.

对于像能量、体积这样的既具有热力学意义又具有纯力学意义的量来讲，涨落的概念是不言而喻的.但是像熵和温度这样的量，它们的确定必须涉及在一段时间内对物体的考察.例如，设 $S(E,V)$ 为物体的平衡熵，它作为（平均）能量和体积的函数.熵的涨落将理解为函数 $S(E,V)$ 的变化，这时 $S(E,V)$ 形式上看成是能量和体积（涨落着的）精确值的函数.

在前几节中我们已经看到，涨落的概率 $w$ 正比于 $\mathrm{e}^{S_{t}}$，其中 $S_{t}$
为闭合系统（即物体作为整体）的总熵.这也可以等效地写成：
\begin{equation*}
  w\propto\mathrm{e}^{\Delta S_{t}},
\end{equation*}
式中 $\Delta S_{t}$ 是熵在涨落时的改变.根据公式(20.8)，我们有：
\begin{equation*}
  \Delta S_{t}=-R_{\min}/T_{0},
\end{equation*}
式中 $R_{\min}$ 是使物体中该小部分以可逆方式发生给定的热力学量改变所必需的最小功（物体的其余部分相对于这一部分来讲起着介质的作用）.因此，
\begin{equation}
  w\propto\mathrm{e}^{-\frac{R_{\min}}{T_{0}}}.
\end{equation}
我们用表达式
\begin{equation*}
  R_{\min}=\Delta E-T_{0}\Delta S+P_{0}\Delta V
\end{equation*}
来代替 $R_{\min}$，式中 $\Delta E,\Delta S,\Delta V$ 分别是物体中该小部分的能量、熵和体积在涨落时的改变，$T_{0}$ 和 $P_{0}$ 是介质的温度和压强，亦即物体温度和压强的平衡（平均）值.以后我们将略去作为涨落量系数的所有各量的角标0；总是指它们的平衡值.因此我们有：
\begin{equation}
  w\propto\exp\Bigl(-\frac{\Delta E-T\Delta S+P\Delta V}{T}\Bigr).
\end{equation}

应当注意：在这种形式下，这个公式适用于任何涨落而不拘大小；至于这里的大涨落，在这种涨落下，例如 $\Delta E$ 可以同物体的这一小部分的能量相当，但比起整个物体的能量自然仍旧很小.应用于小涨落（通常就是这种情形）时，公式(112.2)给出下述结果.

把 $\Delta E$ 展开成级数，我们得到（参看§21）：
\begin{equation*}
  \Delta E-T\Delta S+P\Delta V=\frac{1}{2}\Bigl[
  \frac{\partial^{2}E}{\partial S^{2}}(\Delta S)^{2}
  +2\frac{\partial^{2}E}{\partial S\,\partial V}\Delta S\,\Delta V
  +\frac{\partial^{2}E}{\partial V^{2}}(\Delta V)^{2}\Bigr].
\end{equation*}
很容易看出：这个式子又可以写成形式
\begin{equation*}
  \frac{1}{2}\Bigl[\Delta S\,\Delta\Bigl(\frac{\partial E}{\partial S}\Bigr)_{V}
  +\Delta V\,\Delta\Bigl(\frac{\partial E}{\partial V}\Bigr)_{S}\Bigr]
  =\frac{1}{2}\left(\Delta S\,\Delta T-\Delta P\,\Delta V\right).
\end{equation*}
于是我们得到涨落概率(112.2)的形式如下：
\begin{equation}
  w\propto\exp\frac{\Delta P\,\Delta V-\Delta T\,\Delta S}{2T}.
\end{equation}

从这个普遍公式可以求出各种热力学量的涨落.首先我们选取 $V$ 和 $T$ 作为自变量.于是
\begin{equation*}
  \Delta S=\left(\frac{\partial S}{\partial T}\right)_{V}\Delta T
  +\left(\frac{\partial S}{\partial V}\right)_{T}\Delta V
  =\frac{C_{v}}{T}\Delta T+\left(\frac{\partial P}{\partial T}\right)_{V}\Delta V,
\end{equation*}
\begin{equation*}
  \Delta P=\left(\frac{\partial P}{\partial T}\right)_{V}\Delta T
  +\left(\frac{\partial P}{\partial V}\right)_{T}\Delta V
\end{equation*}
（参看(16.3)）.把这些式子代入(112.3)中的指数，我们就发现：$\Delta V\,\Delta T$ 项相消，而剩下：
\begin{equation}
  w\propto\exp\Bigl[-\frac{C_{v}}{2T^{2}}(\Delta T)^{2}
  +\frac{1}{2T}\Bigl(\frac{\partial P}{\partial V}\Bigr)_{T}(\Delta V)^{2}\Bigr].
\end{equation}
这个式子分解成两个因子，各自只与 $\Delta T$ 或 $\Delta V$ 有关.换句话说，温度和体积的涨落是统计独立的，因而
\begin{equation}
  \left\langle\Delta T\,\Delta V\right\rangle=0.
\end{equation}
把(112.4)中的两个因子逐一与高斯分布的普遍公式(110.6)比较，我们求出温度\footnote{如果 $T$ 用开尔文来计量，则 $\langle(\Delta T)^{2}\rangle=\dfrac{kT^{2}}{C_{v}}$.}和体积的均方涨落的表达式如下：
\begin{equation}
  \left\langle(\Delta T)^{2}\right\rangle=\frac{T^{2}}{C_{v}},
\end{equation}
\begin{equation}
  \left\langle(\Delta V)^{2}\right\rangle
  =-T\left(\frac{\partial V}{\partial P}\right)_{T}.
\end{equation}
这两个量为正，由热力学不等式 $C_{v}>0$ 和
$\left(\dfrac{\partial P}{\partial V}\right)_{T}<0$ 保证.

现在我们选取 $P$ 和 $S$ 作为(112.3)中的自变量.于是
\begin{equation*}
\begin{aligned}
  \Delta V&=\left(\frac{\partial V}{\partial P}\right)_{S}\Delta P
  +\left(\frac{\partial V}{\partial S}\right)_{P}\Delta S,\\
  \Delta T&=\left(\frac{\partial T}{\partial S}\right)_{P}\Delta S
  +\left(\frac{\partial T}{\partial P}\right)_{S}\Delta P
  =\frac{T}{C_{p}}\Delta S+\left(\frac{\partial T}{\partial P}\right)_{S}\Delta P.
\end{aligned}
\end{equation*}
但是根据关系式 $\mathrm{d}W=T\,\mathrm{d}S+V\,\mathrm{d}P$，我们有：
\begin{equation*}
  \left(\frac{\partial V}{\partial S}\right)_{P}
  =\frac{\partial^{2}W}{\partial P\,\partial S}
  =\left(\frac{\partial T}{\partial P}\right)_{S},
\end{equation*}
因此
\begin{equation*}
  \Delta V=\left(\frac{\partial V}{\partial P}\right)_{S}\Delta P
  +\left(\frac{\partial T}{\partial P}\right)_{S}\Delta S.
\end{equation*}
根据公式 $C_{p}=T\left(\dfrac{\partial S}{\partial T}\right)_{P}$，我们有：
\begin{equation*}
  \Delta T=\frac{T}{C_{p}}\Delta S
  +\left(\frac{\partial T}{\partial P}\right)_{S}\Delta P.
\end{equation*}
把 $\Delta V$ 和 $\Delta T$ 代入(112.3)，得：
\begin{equation}
  w\propto\exp\Bigl[\frac{1}{2T}\Bigl(\frac{\partial V}{\partial P}\Bigr)_{S}
  (\Delta P)^{2}-\frac{1}{2C_{p}}(\Delta S)^{2}\Bigr].
\end{equation}
如同(112.4)，这个式子也分解成两个因子，分别与 $\Delta P$ 和 $\Delta S$ 有关.换句话说，熵和压强的涨落是统计独立的\footnote{成对量 $T,V$ 和 $S,P$ 的统计独立性可以由下述的论证直接看出. 如果对给定为§111各式中的 $x_{i}$ 取 $x_{1}=\Delta S$，$x_{2}=\Delta V$，则相应的 $X_{i}$ 为（参看§22）：$X_{1}=\Delta T/T$，$X_{2}=-\Delta P/T$. 但 $i\neq k$ 时 $\langle x_{i}X_{k}\rangle=0$（根据普遍关系式(111.8)），因此得出(112.5)和(112.9).}，因而
\begin{equation}
  \left\langle\Delta S\,\Delta P\right\rangle=0.
\end{equation}
对于熵和压强的均方涨落，我们求出：
\begin{equation}
  \left\langle(\Delta S)^{2}\right\rangle=C_{p},
\end{equation}
\begin{equation}
  \left\langle(\Delta P)^{2}\right\rangle
  =-T\left(\frac{\partial P}{\partial V}\right)_{S}.
\end{equation}

从我们所得到的这些公式可以看出：可加性热力学量（体积和熵）的均方涨落正比于物体中它们所属的那一部分的大小（体积）.与此相应，这些量的均方根涨落正比于体积的平方根，因而相对涨落反比于体积的平方根；这与§2中所作的普遍论断（公式(2.5)）完全一致.但是对于温度和压强这样的量，则均方根涨落本身就反比于体积的平方根.

关系式(112.7)确定了物体中粒子数给定为 $N$ 的某一部分的体积涨落.把这个方程式的两边除以
$N^{2}$，我们就求出单粒子的体积涨落
\begin{equation}
  \left\langle\left[\Delta\Bigl(\frac{V}{N}\Bigr)\right]^{2}\right\rangle
  =-\frac{T}{N^{2}}\left(\frac{\partial V}{\partial P}\right)_{T}.
\end{equation}
显然，这个涨落与我们是在恒定体积下还是在恒定粒子数下考虑涨落必然是无关的.因此由(112.12)可以求出物体给定体积内粒子数的涨落.因为在这种情形下 $V$ 是给定的量，所以应当写成：
\begin{equation*}
  \Delta\Bigl(\frac{V}{N}\Bigr)=V\Delta\Bigl(\frac{1}{N}\Bigr)
  =-\frac{V}{N^{2}}\Delta N.
\end{equation*}
以此代入(112.12)，得：
\begin{equation}
  \left\langle(\Delta N)^{2}\right\rangle
  =-\frac{TN^{2}}{V^{2}}\left(\frac{\partial V}{\partial P}\right)_{T}.
\end{equation}
对于某些计算，这个公式的另一种形式较方便.因为微商
$\left(\dfrac{\partial V}{\partial P}\right)_{T}$ 是在恒定的 $N$ 下取的，我们可写：
\begin{equation*}
  -\frac{N^{2}}{V^{2}}\left(\frac{\partial V}{\partial P}\right)_{T,N}
  =N\left(\frac{\partial}{\partial P}\frac{N}{V}\right)_{T,N}.
\end{equation*}
但是粒子数 $N$ 作为 $P,T,V$ 的函数，由于可加性的要求必须具有 $N=Vf(P,T)$ 的形式（参看§24）；换句话说，$\dfrac{N}{V}$ 只是 $P$ 和 $T$ 的函数，因而在恒定的 $N$ 下或 $V$ 下对
$\dfrac{N}{V}$ 求导都无关紧要，因此我们可写：
\begin{equation*}
\begin{aligned}
  N\left(\frac{\partial}{\partial P}\frac{N}{V}\right)_{T,N}
  &=\frac{N}{V}\left(\frac{\partial N}{\partial P}\right)_{T,V}\\
  &=\left(\frac{\partial N}{\partial P}\right)_{T,V}
  \left(\frac{\partial P}{\partial\mu}\right)_{T,V}
  =\left(\frac{\partial N}{\partial\mu}\right)_{T,V}
\end{aligned}
\end{equation*}
（以上用到了关系式 $\dfrac{N}{V}=\left(\dfrac{\partial P}{\partial\mu}\right)_{T,V}$，它来自式
(24.14) $\mathrm{d}\Omega=V\mathrm{d}P=S\,\mathrm{d}T+N\,\mathrm{d}\mu$）.因此我们得到粒子数涨落的公式如下\footnote{这个公式可以很容易地直接从吉布斯分布得出. 根据平均值的定义，我们有：
\begin{equation*}
  \bar{N}=\mathrm{e}^{\Omega/T}\sum_{N}N\,\mathrm{e}^{\mu N/T}
  \sum_{n}\mathrm{e}^{-E_{n}N/T}.
\end{equation*}
把这个式子对 $\mu$ 微分（在恒定的 $V$ 和 $T$ 下），得
\begin{equation*}
  \frac{\partial\bar{N}}{\partial\mu}
  =\frac{1}{T}\mathrm{e}^{\Omega/T}\sum_{N}
  \Bigl(N^{2}+N\frac{\partial\Omega}{\partial\mu}\Bigr)
  \mathrm{e}^{\mu N/T}\sum_{n}\mathrm{e}^{-E_{n}N/T}
  =\frac{1}{T}\Bigl(\langle N^{2}\rangle
  +\bar{N}\frac{\partial\Omega}{\partial\mu}\Bigr).
\end{equation*}
但是 $\dfrac{\partial\Omega}{\partial\mu}=-\bar{N}$，所以
\begin{equation*}
  \frac{\partial\bar{N}}{\partial\mu}
  =\frac{1}{T}\bigl(\langle N^{2}\rangle-\bar{N}^{2}\bigr)
  =\frac{1}{T}\bigl\langle(\Delta N)^{2}\bigr\rangle,
\end{equation*}
由此得(112.14)式.\par
从吉布斯分布出发，也可以得到其它热力学量涨落的表达式.}：
\begin{equation}
  \left\langle(\Delta N)^{2}\right\rangle
  =T\left(\frac{\partial N}{\partial\mu}\right)_{T,V}.
\end{equation}

除了上面所考虑的各热力学量以外，物体还由它相对于介质的宏观运动动量 $P$ 来描述.在平衡状态没有任何宏观运动，即 $\boldsymbol{P}=0$，但是，运动可由涨落引起；现在我们来确定这种涨落的概率.在这种情形下，最小功 $R_{\min}$ 就等于物体的动能：
\begin{equation*}
  R_{\min}=\frac{P^{2}}{2M}=\frac{Mv^{2}}{2},
\end{equation*}
式中 $M$ 是物体的质量，$\boldsymbol{v}=\dfrac{\boldsymbol{P}}{M}$ 是宏观运动的速度.因此我们得到所求的概率：
\begin{equation}
  w\propto\exp\Bigl(-\frac{Mv^{2}}{2T}\Bigr).
\end{equation}
应当注意：速度涨落是与其它热力学量的涨落统计独立的.速度的每一直角坐标分量的均方涨落都等于
\begin{equation}
  \left\langle(\Delta v_{x})^{2}\right\rangle=\frac{T}{M},
\end{equation}
它与物体的质量成反比.

从所得到的公式可以看出：像能量、体积、压强、速度这些量的经典均方涨落在绝对零度下（正比于温度的一次方）变成零.凡是具有纯力学意义的热力学量都具有这种共同性质，但是一般说来，像熵和温度这样的纯热力学量则没有这种性质.

温度涨落公式(112.6)也可以从另一种观点来解释.我们知道，温度的概念可以通过吉布斯分布引入；这时温度看成是确定这个分布的参量.用于孤立物体时，吉布斯分布可完全描述其统计性质，不精确性仅在于：它导致物体的总能量有虽然很小但毕竟不等于零的涨落，而这种涨落实际上不应当存在（参看§28末尾）.相反地，如果认为物体的能量是一个给定的量，那么就不能赋予物体一个完全确定的温度，而必须认为物体的温度应存在由公式(112.6)所确定的涨落，其中 $C_{v}$ 是物体作为整体的热容.显然这个量表征可以给孤立物体确定温度的精确程度.

\begin{xiti}

\noindent{\heiti 习题1}\quad 试求能量的均方涨落（用 $V$ 和 $T$ 作为自变量）.

解：我们有：
\begin{equation*}
  \Delta E=\left(\frac{\partial E}{\partial V}\right)_{T}\Delta V
  +\left(\frac{\partial E}{\partial T}\right)_{V}\Delta T
  =\Bigl[T\Bigl(\frac{\partial P}{\partial T}\Bigr)_{V}-P\Bigr]\Delta V
  +C_{v}\Delta T.
\end{equation*}
平方后取平均，得
\begin{equation*}
  \left\langle(\Delta E)^{2}\right\rangle
  =-\Bigl[T\Bigl(\frac{\partial P}{\partial T}\Bigr)_{V}-P\Bigr]^{2}
  T\Bigl(\frac{\partial V}{\partial P}\Bigr)_{T}+C_{v}T^{2}.
\end{equation*}

\noindent{\heiti 习题2}\quad 试求 $\left\langle(\Delta W)^{2}\right\rangle$（$P$ 和
$S$ 为变量）.

解：
\begin{equation*}
  \left\langle(\Delta W)^{2}\right\rangle
  =-TV^{2}\left(\frac{\partial P}{\partial V}\right)_{S}+T^{2}C_{p}.
\end{equation*}

\noindent{\heiti 习题3}\quad 试求 $\langle\Delta T\,\Delta P\rangle$（$V$ 和 $T$ 为变量）.

解：
\begin{equation*}
  \left\langle\Delta T\,\Delta P\right\rangle
  =\frac{T^{2}}{C_{v}}\left(\frac{\partial P}{\partial T}\right)_{V}.
\end{equation*}

\noindent{\heiti 习题4}\quad 试求 $\langle\Delta V\,\Delta P\rangle$（$V$ 和 $T$ 为变量）.

解：
\begin{equation*}
  \left\langle\Delta V\,\Delta P\right\rangle=-T.
\end{equation*}

\noindent{\heiti 习题5}\quad 试求 $\langle\Delta S\,\Delta V\rangle$（$V$ 和 $T$ 为变量）.

解：
\begin{equation*}
  \left\langle\Delta S\,\Delta V\right\rangle
  =\left(\frac{\partial V}{\partial T}\right)_{P}T.
\end{equation*}

\noindent{\heiti 习题6}\quad 试求 $\langle\Delta S\,\Delta T\rangle$（$V$ 和 $T$ 为变量）.

解：
\begin{equation*}
  \left\langle\Delta S\,\Delta T\right\rangle=T.
\end{equation*}

\noindent{\heiti 习题7}\quad 试求一垂直悬挂的数学摆的偏角均方涨落.

解：设 $l$ 为摆长，$m$ 为其质量，$\varphi$ 为偏离铅垂线的角度.在这种情形下，最小功
$R_{\min}$ 就是当摆偏离时对重力所作的机械功；小 $\varphi$ 下，
$R_{\min}=\dfrac{1}{2}mg\cdot l\varphi^{2}$.因此
\begin{equation*}
  \left\langle\varphi^{2}\right\rangle=\frac{T}{mgl}.
\end{equation*}

\noindent{\heiti 习题8}\quad 试求一根拉紧的弦上各点偏移的均方涨落.

解：设 $l$ 为弦长，$F$ 为张力.我们考虑弦上与其一端相距 $x$ 处的一点，并设这点的横向位移为
$y$.为了确定 $\langle y^{2}\rangle$，我们必须考虑当 $x$ 这一点有给定位移 $y$ 时弦的平衡形状；它由从弦的两固定端至 $(x,y)$ 点间的两条直线段构成.在弦的这种形变下，所作的功等于
\begin{equation*}
  R_{\min}=F\bigl(\sqrt{x^{2}+y^{2}}-x\bigr)
  +F\bigl[\sqrt{(l-x)^{2}+y^{2}}-(l-x)\bigr]
  \approx\frac{Fy^{2}}{2}\Bigl(\frac{1}{x}+\frac{1}{l-x}\Bigr).
\end{equation*}
由此我们求出均方值为
\begin{equation*}
  \left\langle y^{2}\right\rangle=\frac{T}{Fl}\,x(l-x).
\end{equation*}

\noindent{\heiti 习题9}\quad 试求弦上不同两点位移涨落乘积的平均值.

解：设 $y_{1},y_{2}$ 各为弦上与其一端相距 $x_{1},x_{2}$ 处的两点的横位移（并设
$x_{2}>x_{1}$）.在给定的 $y_{1}$ 和 $y_{2}$ 下，弦的平衡形状由三条直线段构成，所作的功为
\begin{equation*}
  R_{\min}=\frac{F}{2}\Bigl[y_{1}^{2}\frac{x_{2}}{x_{1}(x_{2}-x_{1})}
  +y_{2}^{2}\frac{l-x_{1}}{(l-x_{2})(x_{2}-x_{1})}
  -2y_{1}y_{2}\frac{1}{x_{2}-x_{1}}\Bigr].
\end{equation*}
按照(111.9)，我们求出
\begin{equation*}
  \left\langle y_{1}y_{2}\right\rangle=\frac{T}{Fl}\,x_{1}(l-x_{2}).
\end{equation*}

\end{xiti}

\section{理想气体中的涨落}

把 $V=\dfrac{NT}{P}$ 代入(112.13)，可以对通常理想气体中给定的某一相对较小的体积求出其粒子数的均方涨落.这给出如下的简单结果：
\begin{equation}
  \left\langle(\Delta N)^{2}\right\rangle=N.
\end{equation}
因此粒子数的相对涨落等于粒子平均数平方根的倒数：
\begin{equation*}
  \frac{\sqrt{\left\langle(\Delta N)^{2}\right\rangle}}{N}
  =\frac{1}{\sqrt{N}}.
\end{equation*}

计算玻色或费米理想气体中粒子数的涨落，要用到(112.14)式，其中的 $N$，作为 $\mu,T$ 和
$V$ 的函数，用表达式(56.5)通过积分相应的分布函数得到.我们不在这里写出得到的复杂式子，而只注意下述情况.我们知道，玻色气体在 $T<T_{0}$ 的温度下（参看§62），压强与体积无关；换句话说，它的压缩率变成无限大.根据(112.13)，由此就会得出结论：粒子数的涨落也变成无限大.这意味着：在计算低温下玻色气体中的涨落时，其粒子间的相互作用不论多弱也不能忽略不计；考虑到在任何实际气体中必然存在相互作用，涨落将有限.

我们进一步考虑气体粒子按不同量子态分布的涨落.我们重新考虑粒子的量子态（包括它们平动的不同状态），并设 $n_{k}$ 为态占有数.

我们考虑处于第 $k$ 个量子态的一组 $n_{k}$ 个粒子；由于这一组粒子与气体中其余粒子是完全统计独立的（参看§37），所以对之可用公式(112.14)：
\begin{equation}
  \left\langle(\Delta n_{k})^{2}\right\rangle=T\frac{\partial\bar{n}_{k}}{\partial\mu}.
\end{equation}
上式用于费米气体时，应取
\begin{equation*}
  \bar{n}_{k}=\frac{1}{\mathrm{e}^{(\varepsilon_{k}-\mu)/T}+1}.
\end{equation*}
微分后，得：
\begin{equation}
  \left\langle(\Delta n_{k})^{2}\right\rangle=\bar{n}_{k}(1-\bar{n}_{k}).
\end{equation}
对于玻色气体，用类似方式可得：
\begin{equation}
  \left\langle(\Delta n_{k})^{2}\right\rangle=\bar{n}_{k}(1+\bar{n}_{k}).
\end{equation}
对于玻尔兹曼气体，把 $n_{k}=\exp\Bigl(\dfrac{\mu-\varepsilon_{k}}{T}\Bigr)$ 代入，自然地得到公式
\begin{equation}
  \left\langle(\Delta n_{k})^{2}\right\rangle=\bar{n}_{k},
\end{equation}
这正是(113.3)和(113.4)在 $\bar{n}_{k}\ll1$ 时所趋向的结果.

我们现在把(113.3)和(113.4)对总共包含 $N_{j}=\sum n_{k}$ 个粒子的 $G_{j}$ 个相邻状态求和.由于上述不同 $n_{k}$ 的涨落统计独立性，我们得到：
\begin{equation}
  \left\langle(\Delta N_{j})^{2}\right\rangle
  =G_{j}\bar{n}_{j}(1\mp\bar{n}_{j})
  =\bar{N}_{j}\Bigl(1\mp\frac{\bar{N}_{j}}{G_{j}}\Bigr),
\end{equation}
式中 $\bar{n}_{j}$ 是相邻各 $\bar{n}_{k}$ 的共同值，而 $\bar{N}_{j}=\bar{n}_{j}G_{j}$.
作为特例，以上公式可用于黑体辐射（平衡光子玻色气体），这时，在(113.4)中必须令
$\mu=0$.我们考虑（体积 $V$ 内）频率相近（在小间隔 $\Delta\omega_{j}$ 内）的一组光子量子态；这样量子态的数目等于
\begin{equation*}
  G_{j}=\frac{V\omega_{j}^{2}\Delta\omega_{j}}{\pi^{2}c^{3}}
\end{equation*}
（参看(63.3)）.在这个频率范围内量子的总能量为
\begin{equation*}
  E_{\Delta\omega_{j}}=N_{j}\hbar\omega_{j}.
\end{equation*}
把(113.6)乘以 $(\hbar\omega_{j})^{2}$ 并略去角标 $j$，可得黑体辐射在该频率间隔
$\Delta\omega$ 内能量 $E_{\Delta\omega}$ 涨落的表达式如下：
\begin{equation}
  \left\langle(\Delta E_{\Delta\omega})^{2}\right\rangle
  =\hbar\omega\cdot E_{\Delta\omega}
  +\frac{\pi^{2}c^{3}(E_{\Delta\omega})^{2}}{V\omega^{2}\Delta\omega}.
\end{equation}
这个式子由爱因斯坦（1909）首先推出.

\begin{xiti}

试求电子气在比简并温度低得多的温度下的 $\left\langle(\Delta N)^{2}\right\rangle$.

解：在计算 $\left(\dfrac{\partial N}{\partial\mu}\right)_{T,V}$ 时，我们可以用在绝对零度下
$\mu$ 的表达式(57.3).由简单计算得：
\begin{equation*}
  \left\langle(\Delta N)^{2}\right\rangle
  =\frac{3^{1/3}mT}{\pi^{4/3}\hbar^{2}}
  \left(\frac{N}{V}\right)^{1/3}V.
\end{equation*}

\end{xiti}

\section{泊松公式}

知道了气体给定体积内粒子数的均方涨落(113.1)后，我们就可以对粒子数涨落的概率写出相应的高斯分布：
\begin{equation}
  w(N)\,\mathrm{d}N=\frac{1}{\sqrt{2\pi\bar{N}}}
  \exp\Bigl\{-\frac{(N-\bar{N})^{2}}{2\bar{N}}\Bigr\}\cdot\mathrm{d}N.
\end{equation}
但是这个公式只适用于小涨落. $N-\bar{N}$ 之差必须比粒子数 $\bar{N}$ 本身小得多.

如果气体中所选定的体积 $V$ 足够小，那么其中的粒子数也不大，也值得考虑 $N-\bar{N}$ 同
$\bar{N}$ 可比的大涨落.应当注意：这个问题只对玻尔兹曼气体有意义，因为在费米气体或玻色气体中，这样的涨落的概率只有当体积小到量子涨落占优势的时候才变得显著起来.

解决上述问题的最简单方法如下.设 $V_{0}$ 和 $N_{0}$ 为气体的总体积和其中的粒子总数，并设
$V$ 为比 $V_{0}$ 小得多的一部分体积.由于气体的均匀性，显然某一给定粒子处于体积 $V$ 内的概率就等于比值 $\dfrac{V}{V_{0}}$，而 $N$ 个给定粒子同时处于 $V$ 内的概率为
$\left(\dfrac{V}{V_{0}}\right)^{N}$.同理，一个粒子不在体积 $V$ 内的概率等于
$\dfrac{V_{0}-V}{V_{0}}$，而 $N_{0}-N$ 个给定粒子同时不在的概率为
$\left(1-\dfrac{V}{V_{0}}\right)^{N_{0}-N}$.因此 $N$ 个任意的分子一起处于体积 $V$ 内的概率 $w_{N}$ 由下式给出：
\begin{equation}
  w_{N}=\frac{N_{0}!}{N!\,(N_{0}-N)!}
  \left(\frac{V}{V_{0}}\right)^{N}
  \left(1-\frac{V}{V_{0}}\right)^{N_{0}-N}.
\end{equation}
式中引入了因子 $\dfrac{N_{0}!}{N!\,(N_{0}-N)!}$，它是从 $N_{0}$ 个粒子中选取 $N$ 个的可能方式数.

在我们感兴趣的情形中，$V\ll V_{0}$，而粒子数 $N$ 虽然可能与它的平均值 $\bar{N}$ 相差很大，但比起气体中的粒子总数 $N_{0}$ 来自然还是小得多.因此我们可以令
$\dfrac{N_{0}!}{(N_{0}-N)!}\approx N_{0}^{N}$，并略去幂指数中的 $N$，于是得到：
\begin{equation*}
  w_{N}=\frac{1}{N!}\left(\frac{N_{0}V}{V_{0}}\right)^{N}
  \left(1-\frac{V}{V_{0}}\right)^{N_{0}}.
\end{equation*}
但是 $\dfrac{N_{0}V}{V_{0}}$ 并不是别的，而就是体积 $V$ 内粒子数的平均值 $\bar{N}$.因此我们有：
\begin{equation*}
  w_{N}=\frac{\bar{N}^{N}}{N!}
  \left(1-\frac{\bar{N}}{N_{0}}\right)^{N_{0}}.
\end{equation*}
最后借助于熟知的公式
\begin{equation*}
  \lim_{n\to\infty}\left(1-\frac{x}{n}\right)^{n}=\mathrm{e}^{-x},
\end{equation*}
将 $\left(1-\dfrac{\bar{N}}{N_{0}}\right)^{N_{0}}$ 在大 $N_{0}$ 下用 $\mathrm{e}^{-\bar{N}}$ 代替，最后得所求的概率分布如下式\footnote{对于小涨落（$|N-\bar{N}|\ll\bar{N}$，且 $\bar{N}$ 很大），这个公式自然地过渡到公式(114.1). 看出这一点很容易，只要对大数 $N$ 的阶乘应用斯特林渐近公式
\begin{equation*}
  N!=\sqrt{2\pi N}\cdot N^{N}\mathrm{e}^{-N},
\end{equation*}
并把 $\ln w_{N}$ 按 $N-\bar{N}$ 的幂次展开.}
\begin{equation}
  w_{N}=\frac{\bar{N}^{N}\mathrm{e}^{-\bar{N}}}{N!}.
\end{equation}
这就是所谓泊松公式.容易看出，它满足归一条件
\begin{equation*}
  \sum_{N=0}^{\infty}w_{N}=1.
\end{equation*}
用这个分布可计算粒子数的均方涨落：
\begin{equation*}
\begin{aligned}
  \left\langle N^{2}\right\rangle
  &=\sum_{N=0}^{\infty}N^{2}w_{N}
  =\mathrm{e}^{-\bar{N}}\sum_{N=1}^{\infty}
  \frac{\bar{N}^{N}N}{(N-1)!}=\\
  &=\mathrm{e}^{-\bar{N}}\sum_{N=2}^{\infty}\frac{\bar{N}^{N}}{(N-2)!}
  +\mathrm{e}^{-\bar{N}}\sum_{N=1}^{\infty}\frac{\bar{N}^{N}}{(N-1)!}
  =\bar{N}^{2}+\bar{N}.
\end{aligned}
\end{equation*}
由此，所求的涨落仍有以前的值：
\begin{equation}
  \left\langle(\Delta N)^{2}\right\rangle
  =\left\langle N^{2}\right\rangle-\bar{N}^{2}=\bar{N}.
\end{equation}
因此，不仅对很大的 $\bar{N}$，而且一般对任何 $\bar{N}$ 值，粒子数的均方涨落都等于 $\bar{N}$.

应当注意：公式(114.3)也可以直接从吉布斯分布得出.根据吉布斯分布，气体 $N$ 个粒子同时按不同量子态的分布由下式确定：
\begin{equation*}
  \exp\frac{\Omega+\mu N-\sum\varepsilon_{k}}{T},
\end{equation*}
式中 $\sum\varepsilon_{k}$ 是各个粒子能量之和.为了获得所求的概率 $w_{N}$，必须把这个表达式对分配到这个体积 $V$ 中的粒子所有状态求和.如果对每一个粒子的状态独立求和，得到的结果必须除以 $N$!（参看§41），于是得：
\begin{equation*}
  w_{N}=\frac{\mathrm{e}^{\Omega/T}}{N!}
  \left(\sum_{k}\exp\frac{\mu-\varepsilon_{k}}{T}\right)^{N}
  =\frac{\mathrm{e}^{\Omega/T}}{N!}
  \left(\sum_{k}\bar{n}_{k}\right)^{N}.
\end{equation*}
但是这里的和式不是别的，而就是该体积内粒子数的平均值.因此我们求出：
\begin{equation*}
  w_{N}=\text{常数}\cdot\frac{\bar{N}^{N}}{N!},
\end{equation*}
然后归一化条件给出：常数 $=\mathrm{e}^{-\bar{N}}$\footnote{这就是说，$\Omega=-PV=-\bar{N}T$，与理想气体的物态方程一致.}，重新回到公式(114.3).

\section{溶液中的涨落}

溶液中热力学量的涨落，可以用§112中处理相同粒子所构成物体中的涨落的同样方法计算.如果预先进行下述的考虑，则相应的计算可以大大简化.

我们来考察一小部分溶液，其中包含给定数目 $N$ 的溶剂分子，试图计算这一部分中溶质分子数
$n$ 的平均涨落，也就是其浓度 $c=\dfrac{n}{N}$ 的涨落.为此，我们必须考虑溶液在给定非平衡值
$n$ 下可能的最完全平衡（参看§111第一个脚注）.将浓度取定，不影响溶液的这一小部分建立同其余部分之间涉及彼此能量交换或体积变化的平衡.前者意味着在整个溶液中温度保持常数（参看§9），而后者意味着压强保持常数（参看§12）.因此，要计算均方值
$\left\langle(\Delta c^{2})\right\rangle$，只需考虑在恒温恒压下所发生的浓度涨落.

这个事实本身就意味着：浓度涨落同温度和压强的涨落是统计独立的，换句话说\footnote{更严格的证明可以用§112第三个脚注中所描述的方法给出. 借助于热力学关系 $\mathrm{d}E=T\,\mathrm{d}S-P\,\mathrm{d}V+\mu'\,\mathrm{d}n$（在常数 $N$ 下），(96.1)可改写成形式
\begin{equation*}
  \mathrm{d}R_{\mathrm{min}}=(T-T_{0})\,\mathrm{d}S
  -(P-P_{0})\,\mathrm{d}V+(\mu'-\mu'_{0})\,\mathrm{d}n.
\end{equation*}
由此可以看出：如果选择下列各量作为 $x_{i}$：$x_{1}=\Delta S_{1}$，$x_{2}=\Delta V$，$x_{3}=\Delta n$，那么它们的热力学共轭量为：$X_{1}=\dfrac{\Delta T}{T}$，$X_{2}=-\dfrac{\Delta P}{T}$，$X_{3}=\dfrac{\Delta\mu'}{T}$. 于是由 $\langle x_{3}X_{1}\rangle=0$ 和 $\langle x_{3}X_{2}\rangle=0$ 得出(115.1)式.}，
\begin{equation}
  \langle\Delta T\,\Delta c\rangle=0,\qquad
  \langle\Delta c\,\Delta P\rangle=0.
\end{equation}

根据(96.1)，在恒温恒压下使粒子数 $n$ 改变 $\Delta n$ 所需要的最小功等于
\begin{equation*}
  R_{\min}=\Delta\Phi-\mu'\Delta n,
\end{equation*}
式中 $\mu'$ 是溶质的化学势.把 $\Delta\Phi$ 按 $\Delta n$ 的幂次展开，我们有：
\begin{equation*}
  \Delta\Phi\approx\left(\frac{\partial\Phi}{\partial n}\right)_{P,T}\Delta n
  +\left(\frac{\partial^{2}\Phi}{\partial n^{2}}\right)_{P,T}
  \frac{(\Delta n)^{2}}{2}
  =\mu'\Delta n+\left(\frac{\partial\mu'}{\partial n}\right)_{P,T}
  \frac{(\Delta n)^{2}}{2},
\end{equation*}
于是，
\begin{equation*}
  R_{\min}=\frac{1}{2}\left(\frac{\partial\mu'}{\partial n}\right)_{P,T}
  (\Delta n)^{2}.
\end{equation*}
把这个表达式代入普遍公式(112.1)，并与高斯分布公式(110.5)比较，我们得到所求的粒子数 $n$ 的均方涨落：
\begin{equation}
  \left\langle(\Delta n)^{2}\right\rangle
  =\frac{T}{(\partial\mu'/\partial n)_{P,T}};
\end{equation}
除以 $N^{2}$，得浓度的均方涨落：
\begin{equation}
  \left\langle(\Delta c)^{2}\right\rangle
  =\frac{T}{N(\partial\mu'/\partial c)_{P,T}}.
\end{equation}
后者反比于溶液的这一小部分的物质量（$N$），这原是理所当然的（参看(112.11)后面的讨论）.

对于稀溶液，$\dfrac{\partial\mu'}{\partial n}=\dfrac{T}{n}$，(115.2)式给出：
\begin{equation}
  \left\langle(\Delta n)^{2}\right\rangle=n.
\end{equation}
应当注意，正如所料，这与理想气体中的粒子数涨落公式(113.1)完全类似.

\section{密度涨落的空间关联}

在均匀的各向同性的介质（气体或液体）中，粒子在空间的所有位置都是等概率的，这一论断，在其余所有的粒子可以占据任意位置的条件下，适用于每一个给定的粒子.自然，这个论断并不与下述事实矛盾：由于粒子的相互作用，在不同粒子的相对位置之间必然存在某种关联.譬如说，如果我们同时考虑两个粒子，那么当给定一个粒子的位置时，另一个粒子的不同位置并不等概率.

用 $n(\boldsymbol{r})$ 表示精确（涨落着）的粒子数密度；乘积 $n\,\mathrm{d}V$ 是（在给定的时刻）处于体积元 $\mathrm{d}V$ 内的粒子数.为表征粒子在空间两点位置间的关联，我们引入密度涨落的空间关联函数：
\begin{equation}
  \left\langle\Delta n_{1}\Delta n_{2}\right\rangle
  =\overline{n_{1}n_{2}}-\bar{n}^{2},
\end{equation}
式中 $\Delta n=n-\bar{n}$，而指标1和2区别在空间两点 $\boldsymbol{r}_{1}$ 和
$\boldsymbol{r}_{2}$ 的 $n(\boldsymbol{r})$ 值.在各向同性均匀介质中关联函数仅仅与两点之间距离的绝对值 $r=|\boldsymbol{r}_{2}-\boldsymbol{r}_{1}|$ 有关.当
$\boldsymbol{r}\to\infty$ 时，在点 $\boldsymbol{r}_{1}$ 和 $\boldsymbol{r}_{2}$ 的涨落成为统计独立的，因此关联函数趋于0.

对于用上述方式引入的关联函数，通过以下讨论来解释其意义是有益的.

由于体积 $\mathrm{d}V$ 无限小，其中不可能同时有一个以上的粒子；在其中同时发现两个粒子的概率是高阶的无穷小.所以粒子的平均数 $\bar{n}\,\mathrm{d}V$ 同时也是粒子处于体积元
$\mathrm{d}V$ 中的概率.我们再用 $\bar{n}w_{12}(r)\,\mathrm{d}V_{2}$ 表示在一个粒子处于体积元 $\mathrm{d}V_{1}$ 的条件下，粒子处于体积元 $\mathrm{d}V_{2}$ 的概率
（$r\to\infty$ 时，$w_{12}\to1$）.如上所述，显然平均值
\begin{equation*}
  \left\langle n_{1}\mathrm{d}V_{1}\cdot n_{2}\mathrm{d}V_{2}\right\rangle
  =\bar{n}\,\mathrm{d}V_{1}\cdot\bar{n}w_{12}\,\mathrm{d}V_{2}.
\end{equation*}
因此 $\left\langle n_{1}n_{2}\right\rangle=w_{12}\bar{n}^{2}$.这个等式在
$\boldsymbol{r}_{1}\neq\boldsymbol{r}_{2}$ 时成立，但不能过渡到
$\boldsymbol{r}_{2}\to\boldsymbol{r}_{1}$ 的极限情形，因为推导中没有考虑到，如果点1和点2相同，则处于 $\mathrm{d}V_{1}$ 中的粒子也处于 $\mathrm{d}V_{2}$ 中.容易看出，考虑了这种情况的关系式应当具有形式
\begin{equation}
  \left\langle n_{1}n_{2}\right\rangle
  =\bar{n}^{2}w_{12}+\bar{n}\,\delta(\boldsymbol{r}_{2}-\boldsymbol{r}_{1}).
\end{equation}
实际上，如果我们取一个小体积 $\Delta V$，把(116.2)乘以 $\mathrm{d}V_{1}\mathrm{d}V_{2}$，并遍及这个小体积积分.这时 $\bar{n}^{2}w_{12}$ 项给出一个二阶小量（正比于
$(\Delta V)^{2}$）；而带 $\delta$ 函数的项给出一阶量 $\bar{n}\Delta V$，也本该如此，因为准确到一阶量时只可能有0或1个粒子处于小体积内.

关联函数(116.1)也可适当地分出带 $\delta$ 函数的项，写成形式
\begin{equation}
  \left\langle\Delta n_{1}\Delta n_{2}\right\rangle
  =\bar{n}\,\delta(\boldsymbol{r}_{2}-\boldsymbol{r}_{1})+\bar{n}\,\nu(r),
\end{equation}
式中
\begin{equation}
  \nu(r)=\bar{n}\left[w_{12}(r)-1\right].
\end{equation}
我们将原始量 $\left\langle\Delta n_{1}\Delta n_{2}\right\rangle$ 和函数 $\nu(r)$\footnote{函数 $\nu(r)$ 与§79中引入的函数 $\omega_{12}(r)$ 区别在于归一化：$\bar{n}\omega_{12}=\nu$.}
二者都称为 $\Delta n$ 的关联函数.

现在取某一有限体积 $V$，把等式(116.3)对 $\mathrm{d}V_{1}\mathrm{d}V_{2}$ 积分.引入在这个体积中的粒子总数 $N$（因此 $\bar{n}V=\bar{N}$），得：
\begin{equation*}
  \left\langle(\Delta N)^{2}\right\rangle
  =\bar{N}+\bar{n}\iint\nu(r)\,\mathrm{d}V_{1}\mathrm{d}V_{2},
\end{equation*}
或者把对 $\mathrm{d}V_{1}\mathrm{d}V_{2}$ 的积分转为对其中一个粒子的坐标以及相对坐标
$\boldsymbol{r}=\boldsymbol{r}_{2}-\boldsymbol{r}_{1}$ 的积分，
\begin{equation}
  \int\nu\,\mathrm{d}V=\frac{\left\langle(\Delta N)^{2}\right\rangle}{\bar{N}}-1.
\end{equation}
于是，关联函数对某个体积的积分可用这个体积内总粒子数的均方涨落表示.利用后者的热力学公式(112.13)，可借助热力学量表示这个积分：
\begin{equation}
  \int\nu\,\mathrm{d}V=-\frac{TN}{V^{2}}
  \left(\frac{\partial V}{\partial P}\right)_{T}-1.
\end{equation}
对于普通的（经典的）理想气体，理所当然积分变成零：在这种气体中，不同粒子位置间没有任何关联，因为它们之间没有任何相互作用——既没有直接的，也没有（量子的理想气体中）交换的相互作用.

相反，在液体中（远离临界点），由于液体的压缩率很小，表达式(116.6)中的第一项比1小得多，积分接近于$-1$\footnote{值 $-1$ 似乎对应于液体粒子的彼此不可穿透性，它们被看成是密实的小球.}.液体粒子之间的主要相互作用力，作用半径具有分子线度 $a$ 的数量级.考虑到这些力，关联函数 $\nu(r)$ 将随距离指数衰减，而指数为$\sim-\dfrac{r}{a}$\footnote{但是也存在长程弱相互作用力（范德瓦尔斯力）. 这些力导致在相关函数中出现随距离（按幂律）缓慢衰减的项（参看本教程第九卷）.}.

因为密度涨落和温度涨落是统计独立的，所以在研究密度涨落时，可以把温度当成是不变的.根据定义，物体的总体积也不变.在这些条件下，为使物体偏离平衡所需做的最小功等于物体总自由能的改变 $\Delta F_{t}$.因此涨落的概率
\begin{equation}
  w\propto\exp\Bigl(-\frac{\Delta F_{t}}{T}\Bigr).
\end{equation}
与密度涨落有关的改变 $\Delta F_{t}$ 可以表示为如下形式：
\begin{equation}
  \Delta F_{t}=\frac{1}{2}\iint\varphi(\boldsymbol{r})
  \Delta n_{1}\Delta n_{2}\,\mathrm{d}V_{1}\mathrm{d}V_{2}.
\end{equation}
我们将演示如何用函数 $\varphi(\boldsymbol{r})$ 求出相关函数 $\nu(r)$\footnote{用数学术语说，$\varphi(\boldsymbol{r})$ 是 $\Delta F_{t}$ 对 $n(\boldsymbol{r})$ 的二阶变分导数.}.

我们考虑体积 $V$ 很大但有限的物体，把 $\Delta n$ 展开成傅里叶级数：
\begin{equation}
  \Delta n=\sum_{k}\Delta n_{k}\,\mathrm{e}^{\mathrm{i}\boldsymbol{k}\cdot\boldsymbol{r}},
  \qquad
  \Delta n_{k}=\frac{1}{V}\int\Delta n\,
  \mathrm{e}^{-\mathrm{i}\boldsymbol{k}\cdot\boldsymbol{r}}\mathrm{d}V
\end{equation}
（由于 $\Delta n$ 为实数 $\Delta n_{-k}=\Delta n^{*}_{k}$）.把这些表达式代入(116.8)并积分，所有带乘积
$\Delta n_{k}\Delta n_{k'}\mathrm{e}^{\mathrm{i}(\boldsymbol{k}+\boldsymbol{k}')\cdot\boldsymbol{r}}$
的项，$\boldsymbol{k}'\neq-\boldsymbol{k}$ 时变为零，结果求得
\begin{equation}
  \Delta F_{t}=\frac{V}{2}\sum_{k}|\Delta n_{k}|^{2}\varphi(\boldsymbol{k}),
\end{equation}
式中采用以新宗量 $\boldsymbol{k}$ 标识的同样字母 $\varphi$ 表示函数 $\varphi(\boldsymbol{r})$ 的傅里叶积分分量
\begin{equation}
  \varphi(\boldsymbol{k})=\int\varphi(\boldsymbol{r})\,
  \mathrm{e}^{-\mathrm{i}\boldsymbol{k}\cdot\boldsymbol{r}}\mathrm{d}V.
\end{equation}
因为和式(116.10)的每一项只与 $\Delta n_{k}$ 中的一项有关，所以不同 $\Delta n_{k}$ 的涨落统计独立.每一个平方项 $|\Delta n_{k}|^{2}$ 在和式中出现两次（$\pm\boldsymbol{k}$），因此它的涨落的概率分布由下式给出
\begin{equation*}
  w\propto\exp\Bigl(-\frac{V}{T}\varphi(\boldsymbol{k})
  \left|\Delta n_{k}\right|^{2}\Bigr).
\end{equation*}
最后，考虑到 $|\Delta n_{k}|^{2}$ 是两个独立量的平方和（$\Delta n_{k}$ 是复的），由此求得均方涨落
\begin{equation}
  \left\langle|\Delta n_{k}|^{2}\right\rangle=\frac{T}{V\varphi(\boldsymbol{k})}.
\end{equation}
另一方面，在等式(116.3)的两边乘以
$\exp(-\mathrm{i}\boldsymbol{k}\cdot\boldsymbol{r})
=\exp[-\mathrm{i}\boldsymbol{k}\cdot(\boldsymbol{r}_{2}-\boldsymbol{r}_{1})]$
把它再对 $\mathrm{d}V_{1}\mathrm{d}V_{2}$ 积分，我们得到
\begin{equation}
  \left\langle|\Delta n_{k}|^{2}\right\rangle
  =\frac{\bar{n}}{V}\left[1+\nu(k)\right],\qquad
  \nu(k)=\int\nu(r)\,\mathrm{e}^{\mathrm{i}\boldsymbol{k}\cdot\boldsymbol{r}}\mathrm{d}V.
\end{equation}
最后，把(116.12)代入这里，得所求结果
\begin{equation}
  \nu(k)=\frac{T}{\bar{n}\varphi(k)}-1.
\end{equation}

\begin{xiti}

确定稀薄气体关联函数按 $\dfrac{N}{V}$ 的幂次展开的一次项.

解：从公式(79.2)出发，在一级近似下可以认为，给定的两个粒子外的所有其余粒子相距很远，它们的相互作用可以忽略不计，因此积分后给出 $V^{N-2}$.以同样的精度可以设
$F=F_{\mathrm{id}}$.结果求得
\begin{equation*}
  \nu(r)=\bar{n}\left[\mathrm{e}^{-U(r)/T}-1\right],
\end{equation*}
式中 $U(r)$ 是气体两个粒子的相互作用能.

值得注意，把这个表达式代入(79.1)可给出气体的能量
\begin{equation*}
  E=E_{\mathrm{id}}+\frac{N^{2}}{2V^{2}}V\int U\Bigl(1+\frac{\nu}{\bar{n}}\Bigr)
  \mathrm{d}V=E_{\mathrm{id}}+\frac{N^{2}}{2V}\int U\,\mathrm{e}^{-U/T}\,\mathrm{d}V.
\end{equation*}
当然，这个结果与近理想气体自由能的公式(74.4)，(74.5)一致.

\end{xiti}

\section{简并气体中的密度涨落关联}

正如上一节指出的，在经典理想气体中不同粒子的位置之间通常没有任何关联.但是，在量子力学中根据波函数的对称性原理\footnote{对于费米气体中涨落的关联，В. С. Фурсов 曾研究过（1937），而对玻色气体 А. Я. Галанин 曾研究过（1940）.}，由于理想气体粒子的间接相互作用产生了关联.

简并气体关联函数如何确定的问题，用二次量子化方法（已在§80中计算电子气能量时用过）解决最为简单.

众所周知，在这种方法里粒子数密度对应于算符
\begin{equation*}
  \hat{n}(\boldsymbol{r})=\hat{\psi}^{+}(\boldsymbol{r})\hat{\psi}(\boldsymbol{r});
\end{equation*}
把 $\psi$ 算符(80.5)代入后，用和式表示
\begin{equation}
  \hat{n}(\boldsymbol{r})=\sum_{\sigma\sigma'pp'}
  \hat{a}^{+}_{p\sigma}\hat{a}_{p'\sigma'}
  \psi^{*}_{p\sigma}(\boldsymbol{r})\psi_{p'\sigma'}(\boldsymbol{r}),
\end{equation}
式中求和遍及所有的动量值 $\boldsymbol{p},\boldsymbol{p}'$（对于体积 $V$ 中的自由粒子）以及自旋分量 $\sigma,\sigma'$\footnote{应该提醒：具有自旋的粒子的波函数是旋量，因而(117.1)中波函数的乘积实际上是协变旋量和逆变旋量按旋量指标（不应将之与表示状态自旋本征值分量的指标 $\sigma,\sigma'$ 相混）适当求和后的“标量积”.}.但是由于 $\sigma$ 值不同的自旋波函数有正交性，和式中只有
$\sigma=\sigma'$ 的项实际上不为零.在乘积 $\psi^{*}_{p\sigma}\psi_{p'\sigma}$
中归一化的自旋因子给出1，因此，波函数可以简单地写成坐标平面波的形式
\begin{equation}
  \psi_{p}=\frac{1}{\sqrt{V}}\,
  \mathrm{e}^{\mathrm{i}\boldsymbol{p}\cdot\boldsymbol{r}/\hbar}.
\end{equation}
容易看出，和式(117.1)的所有对角项（$\boldsymbol{p}=\boldsymbol{p}'$）正好给出平均密度
$\bar{n}$：因为算符 $\hat{a}^{+}_{p\sigma}\hat{a}_{p\sigma}$ 就是在给定量子状态下的粒子数
$n_{p\sigma}$，所以这些项的和等于
\begin{equation*}
  \frac{1}{V}\sum_{\sigma p}n_{p\sigma}=\frac{N}{V}=\bar{n}.
\end{equation*}
因此可写
\begin{equation}
  \Delta\hat{n}=\hat{n}(\boldsymbol{r})-\bar{n}
  =\sum'_{\sigma pp'}\hat{a}^{+}_{p\sigma}\hat{a}_{p'\sigma}
  \psi^{*}_{p}\psi_{p'},
\end{equation}
此处求和号上的撇号表示不计其中的对角项.借助这个式子不难计算待求的平均值
$\left\langle\Delta n_{1}\Delta n_{2}\right\rangle$.

平均值的计算分两步.首先应当对粒子状态求量子力学平均.这种平均归结为对给定物理量取相应的对角矩阵元.把属于不同点 $\boldsymbol{r}_{1}$ 和 $\boldsymbol{r}_{2}$ 的两个算符(117.3)式相乘，我们得到一个和式，它包含每次取四个算符
$\hat{a}_{p\sigma},\hat{a}^{+}_{p\sigma}$ 所得的不同乘积项.但是在所有这些乘积中，具有对角矩阵元的只是那些含有两对带等同指标算符 $\hat{a}_{p\sigma},\hat{a}^{+}_{p\sigma}$ 之项，即
\begin{equation*}
  \sum'_{\sigma pp'}\hat{a}^{+}_{p\sigma}\hat{a}_{p'\sigma}
  \hat{a}^{+}_{p'\sigma}\hat{a}_{p\sigma}
  \psi^{*}_{p}(\boldsymbol{r}_{1})\psi_{p'}(\boldsymbol{r}_{1})
  \psi^{*}_{p'}(\boldsymbol{r}_{2})\psi_{p}(\boldsymbol{r}_{2}).
\end{equation*}
这些项构成对角矩阵，有
\begin{equation*}
  \hat{a}_{p'\sigma}\hat{a}^{+}_{p'\sigma}=1\mp n_{p'\sigma},\qquad
  \hat{a}^{+}_{p\sigma}\hat{a}_{p\sigma}=n_{p\sigma}
\end{equation*}
（从此往后，上号记费米统计，而下号记玻色统计）.再把(117.2)式的 $\psi_{p}$ 函数代入，得
\begin{equation*}
  \frac{1}{V^{2}}\sum'_{\sigma pp'}(1\mp n_{p'\sigma})n_{p\sigma}
  \exp\Bigl[\frac{\mathrm{i}}{\hbar}
  (\boldsymbol{p}-\boldsymbol{p}')\cdot(\boldsymbol{r}_{2}-\boldsymbol{r}_{1})\Bigr].
\end{equation*}
这个表达式现在应该进行在统计意义上（即按照粒子在各量子态上的平衡分布）的平均.因为处于不同量子态的粒子彼此行为完全独立，所以粒子数 $n_{p\sigma}$ 和 $n_{p'\sigma}$ 可独立求平均.结果我们得到所求的平均值
\begin{equation}
  \left\langle\Delta n_{1}\Delta n_{2}\right\rangle
  =\frac{1}{V^{2}}\sum'_{\sigma pp'}(1\mp\bar{n}_{p'\sigma})\bar{n}_{p\sigma}
  \exp\Bigl[\frac{\mathrm{i}}{\hbar}
  (\boldsymbol{p}-\boldsymbol{p}')\cdot(\boldsymbol{r}_{2}-\boldsymbol{r}_{1})\Bigr].
\end{equation}

现在照一般方式将对 $\boldsymbol{p},\boldsymbol{p}'$ 的求和化为对
$\dfrac{V\mathrm{d}^{3}p\,V\mathrm{d}^{3}p'}{(2\pi\hbar)^{6}}$ 的积分（这时
$\boldsymbol{p}\neq\boldsymbol{p}'$ 的限制变得不再重要）.积分分为两部分，其中第一部分为
\begin{equation*}
  \sum_{\sigma}\int\bar{n}_{p\sigma}
  \exp\Bigl[\frac{\mathrm{i}}{\hbar}
  (\boldsymbol{p}-\boldsymbol{p}')\cdot(\boldsymbol{r}_{2}-\boldsymbol{r}_{1})\Bigr]
  \frac{\mathrm{d}^{3}p\,\mathrm{d}^{3}p'}{(2\pi\hbar)^{6}}.
\end{equation*}
对 $\dfrac{\mathrm{d}^{3}p'}{(2\pi\hbar)^{3}}$ 的积分给出 $\delta$ 函数
$\delta(\boldsymbol{r}_{2}-\boldsymbol{r}_{1})$，这使得可对被积式其余部分设
$\boldsymbol{r}_{2}-\boldsymbol{r}_{1}=0$；此后有
\begin{equation*}
  \delta(\boldsymbol{r}_{2}-\boldsymbol{r}_{1})\sum_{\sigma}\int\bar{n}_{p\sigma}
  \frac{\mathrm{d}^{3}p}{(2\pi\hbar)^{3}}
  =\bar{n}\,\delta(\boldsymbol{r}_{2}-\boldsymbol{r}_{1}).
\end{equation*}
这正好是公式(116.3)中的第一项.因此，关联函数（(116.3)中的第二项）有下面的表达式：
\begin{equation}
  \nu(r)=\mp\frac{1}{\bar{n}}\sum_{\sigma}
  \left|\int\mathrm{e}^{\mathrm{i}\boldsymbol{p}\cdot\boldsymbol{r}/\hbar}
  \bar{n}_{p\sigma}\frac{\mathrm{d}^{3}p}{(2\pi\hbar)^{3}}\right|^{2}.
\end{equation}
在平衡的气体中，粒子按量子态的分布依从费米或玻色分布公式
\begin{equation}
  \bar{n}_{p\sigma}\equiv\bar{n}_{p}
  =\frac{1}{\mathrm{e}^{(\varepsilon-\mu)/T}\pm1}.
\end{equation}
这些数与 $\sigma$ 无关；因此(117.5)中的对 $\sigma$ 求和就给出因子 $g=2s+1$（$s$
是粒子的自旋）.所以，最后我们得到关联函数的如下公式\footnote{对于玻色气体，这个公式只适于温度比玻色−爱因斯坦凝聚点高的情况（参看习题4）.}
\begin{equation}
  \nu(r)=\mp\frac{g}{\bar{n}}
  \left|\int\frac{\mathrm{e}^{\mathrm{i}\boldsymbol{p}\cdot\boldsymbol{r}/\hbar}}
  {\mathrm{e}^{(\varepsilon-\mu)/T}\pm1}
  \frac{\mathrm{d}^{3}p}{(2\pi\hbar)^{3}}\right|^{2},
\end{equation}
或者，对 $\boldsymbol{p}$ 的取向积分后得
\begin{equation}
  \nu(r)=\mp\frac{g}{4\pi^{4}\bar{n}r^{2}\hbar^{4}}
  \left|\int_{0}^{\infty}
  \frac{\sin(pr/\hbar)\,p\,\mathrm{d}p}
  {\mathrm{e}^{(\varepsilon-\mu)/T}\pm1}\right|^{2}.
\end{equation}
密度涨落傅里叶分量均方值的公式，只要将(117.7)式 $\nu(r)$ 代入普遍公式(116.13)并对坐标积分，也很容易得到\footnote{谨防将气体密度涨落的傅里叶分量 $\Delta n_{k}$ 与粒子的量子态占有数 $\bar{n}_{p}$ 相混.}：
\begin{equation}
  \left\langle|\Delta n_{k}|^{2}\right\rangle
  =\frac{g}{V}\int\bar{n}_{p}(1\mp\bar{n}_{p+\hbar k})
  \frac{\mathrm{d}^{3}p}{(2\pi\hbar)^{3}}.
\end{equation}

从(117.7)式首先可以看出，对于费米气体 $\nu(r)<0$，而对于玻色气体 $\nu(r)>0$.
换而言之，在玻色气体中，一个粒子在空间某点的存在增加了在这点附近发现另一个粒子的概率，即粒子经受独特的吸引力.而在费米气体中，恰好相反，粒子呈现出的类似于排斥力（参看§56末尾的说明）.

根据在本节开头所述，在经典极限下关联函数变为零：在 $\hbar\to0$ 时，(117.7)中被积表达式振荡因子
$\exp\Bigl(\dfrac{\mathrm{i}}{\hbar}\boldsymbol{p}\cdot\boldsymbol{r}\Bigr)$
的频率无限制地增加，积分趋于零.

当 $r\to0$ 时，函数 $\nu(r)$ 趋于常数极限：
\begin{equation}
  \nu(0)=\mp\frac{g}{\bar{n}}
  \left|\int\bar{n}_{p}\frac{\mathrm{d}^{3}p}{(2\pi\hbar)^{3}}\right|^{2}
  =\mp\frac{\bar{n}}{g}.
\end{equation}
我们把(117.8)式应用于 $T=0$ 时的费米气体.在这种情况下，分布函数是阶跃函数：当
$\boldsymbol{p}<\boldsymbol{p}_{\mathrm F}$ 时，$\bar{n}_{p}=1$；当
$\boldsymbol{p}>\boldsymbol{p}_{\mathrm F}$ 时，$\bar{n}_{p}=0$，式中
$p_{\mathrm F}=\hbar(6\pi^{2}\bar{n}/g)^{1/3}$ 是限界动量.因此求得
\begin{equation*}
  \nu(r)=-\frac{g}{4\pi^{4}\hbar^{4}\bar{n}r^{2}}
  \left|\int_{0}^{p_{\mathrm F}}p\sin\frac{pr}{\hbar}\,\mathrm{d}p\right|^{2}.
\end{equation*}
我们考虑不太小的距离，即认为 $\dfrac{p_{\mathrm F}r}{\hbar}\gg1$.与此相应，只保留积分中
$\dfrac{1}{r}$ 的最低次项：
\begin{equation}
  \nu(r)=-\frac{3\hbar}{2\pi^{2}p_{\mathrm F}r^{4}}
  \cos^{2}\frac{p_{\mathrm F}r}{\hbar}
  =-\frac{3\hbar}{4\pi^{2}p_{\mathrm F}r^{4}}
  \Bigl(1+\cos\frac{2p_{\mathrm F}r}{\hbar}\Bigr).
\end{equation}
在比所考虑距离小得多的区间 $\Delta r$ 上，余弦变化很快.对这样的区间求平均，得到
\begin{equation}
  \nu(r)=-\frac{3\hbar}{4\pi^{2}p_{\mathrm F}r^{4}}.
\end{equation}

\begin{xiti}

\noindent{\heiti 习题1}\quad 确定费米气体在 $T=0$ 时密度涨落的低波数
$\bigl(k\ll p_{\mathrm F}/\hbar\bigr)$ 傅里叶分量的均方值.

解：(117.9)中的被积式只有在 $\bar{n}_{p}=1$，$\bar{n}_{p+\hbar k}=0$ 的各点才不等于0（而等于1），这些点在以原点为中心以 $p_{\mathrm F}$ 为半径的球内，同时在中心移了
$\hbar k$ 而半径相同的球以外.在 $\hbar k\ll p_{\mathrm F}$ 下求出这个区域的体积，得到
\begin{equation*}
  \left\langle|\Delta n_{k}|^{2}\right\rangle
  =\frac{\pi gkp^{2}_{\mathrm F}}{(2\pi)^{3}\hbar^{2}V}
  =\frac{3k\hbar}{4p_{\mathrm F}}\,\frac{\bar{n}}{V}.
\end{equation*}

\noindent{\heiti 习题2}\quad 确定费米气体在比简并温度低得多的温度下的关联函数.

解：在(117.8)的积分中设 $\mu\approx\varepsilon_{\mathrm F}=\dfrac{p^{2}_{\mathrm F}}{2m}$，并作如下变换：
\begin{equation*}
  I=\int_{0}^{\infty}
  \frac{p\sin(pr/\hbar)\,\mathrm{d}p}
  {\mathrm{e}^{(\varepsilon-\varepsilon_{\mathrm F})/T}+1}
  =-\hbar\frac{\partial}{\partial r}\int_{0}^{\infty}
  \frac{\cos(pr/\hbar)\,\mathrm{d}p}
  {\mathrm{e}^{(\varepsilon-\varepsilon_{\mathrm F})/T}+1}.
\end{equation*}
分部积分，之后引入新变量 $x=\dfrac{p_{\mathrm F}(p-p_{\mathrm F})}{mT}$.由于 $T$
很小，被积式随着 $|x|$ 的增加而很快地减小，对 $\mathrm{d}x$ 的积分可以扩展为从 $-\infty$ 到
$+\infty$
\begin{equation*}
\begin{aligned}
  I&=-\hbar^{2}\frac{\partial}{\partial r}\frac{1}{r}
  \int_{-\infty}^{\infty}
  \sin\Bigl(\frac{p_{\mathrm F}}{\hbar}r+\lambda xr\Bigr)
  \frac{\mathrm{d}x}{(\mathrm{e}^{x}+1)(\mathrm{e}^{-x}+1)}=\\
  &=-\hbar^{2}\frac{\partial}{\partial r}\biggl[
  \frac{\sin(p_{\mathrm F}r/\hbar)}{r}
  \int_{-\infty}^{\infty}\mathrm{e}^{\mathrm{i}\lambda rx}
  \frac{\mathrm{d}x}{(\mathrm{e}^{x}+1)(\mathrm{e}^{-x}+1)}\biggr]
\end{aligned}
\end{equation*}
（式中 $\lambda=\dfrac{mT}{\hbar p_{\mathrm F}}$）.所得到的积分用
$(\mathrm{e}^{x}+1)^{-1}=u$ 代入化为欧拉B积分，结果得到
\begin{equation*}
  I=\hbar^{2}\frac{\partial}{\partial r}\biggl[
  \frac{\pi\lambda}{\sinh(\pi\lambda r)}
  \sin\frac{p_{\mathrm F}r}{\hbar}\biggr].
\end{equation*}
对于 $r\gg\dfrac{\hbar}{p_{\mathrm F}}$ 的距离，对变化很快的余弦平方取平均，最后得到
\begin{equation*}
  \nu(r)=-\frac{3(mT)^{2}}{4\hbar p^{3}_{\mathrm F}r^{2}}
  \sinh^{-2}\Bigl(\frac{\pi mTr}{\hbar p_{\mathrm F}}\Bigr).
\end{equation*}
在 $T\to0$ 时这个公式变为(117.11).在 $\dfrac{p_{\mathrm F}r}{\hbar}$ 不仅比1大，而且比
$\dfrac{\varepsilon_{\mathrm F}}{T}$ 也大的渐近区内，我们有
\begin{equation*}
  \nu(r)=-\frac{3(mT)^{2}}{\hbar p^{3}_{\mathrm F}r^{2}}
  \exp\Bigl(-\frac{2\pi mTr}{\hbar p_{\mathrm F}}\Bigr).
\end{equation*}

\noindent{\heiti 习题3}\quad 对于玻色气体，在稍稍高于玻色-爱因斯坦凝聚起始点 $T_{0}$
的温度下，求其大距离 $\bigl(r\gg\hbar/\sqrt{mT}\,\bigr)$ 的关联函数.

解：在 $T_{0}$ 点附近化学势 $|\mu|$ 很小（参看§62习题）.这时(117.7)中的积分（记作
$I$）取决于小 $p$ 值区：$\dfrac{\varepsilon}{T}\sim\dfrac{p^{2}}{mT}
\sim\dfrac{|\mu|}{T}\ll1$.因此，被积表达式按 $\varepsilon$ 和 $\mu$ 展开，求得\footnote{此处用了傅里叶变换公式
\begin{equation*}
  \int\frac{\mathrm{e}^{-\varkappa r}}{r}\,
  \mathrm{e}^{-\mathrm{i}\boldsymbol{k}\cdot\boldsymbol{r}}\,\mathrm{d}V
  =\frac{4\pi}{\varkappa^{2}+k^{2}},\qquad
  \int\frac{\mathrm{e}^{\mathrm{i}\boldsymbol{k}\cdot\boldsymbol{r}}}
  {\varkappa^{2}+k^{2}}\,\frac{\mathrm{d}^{3}k}{(2\pi)^{3}}
  =\frac{\mathrm{e}^{-\varkappa r}}{4\pi r}.
\end{equation*}
得出该式的最简单方法如下. 注意到函数 $\varphi=\mathrm{e}^{-\varkappa r}/r$ 满足微分方程
\begin{equation*}
  \Delta\varphi-\varkappa^{2}\varphi=-4\pi\delta(\boldsymbol{r}),
\end{equation*}
方程两边乘以 $\mathrm{e}^{-\mathrm{i}\boldsymbol{k}\cdot\boldsymbol{r}}$ 并在整个空间积分（并且对 $\mathrm{e}^{-\mathrm{i}\boldsymbol{k}\cdot\boldsymbol{r}}\Delta\varphi$ 作两次分部积分），就得到所要的结果.}
\begin{equation*}
  I\approx T\int
  \frac{\mathrm{e}^{\mathrm{i}\boldsymbol{p}\cdot\boldsymbol{r}/\hbar}\,
  \mathrm{d}^{3}p}{(2\pi\hbar)^{3}\left(p^{2}/2m+|\mu|\right)}
  =\frac{mT}{2\pi\hbar^{2}r}
  \exp\Bigl[-r\frac{(2m|\mu|)^{1/2}}{\hbar}\Bigr].
\end{equation*}
最后，
\begin{equation*}
  \nu(r)=\frac{gm^{2}T^{2}}{4\pi^{2}\bar{n}\hbar^{4}r^{2}}
  \exp\biggl[-r\frac{2(2m|\mu|)^{1/2}}{\hbar}\biggr].
\end{equation*}

\noindent{\heiti 习题4}\quad 求玻色气体在 $T<T_{0}$ 时的相关函数.

解：在 $T<T_{0}$ 下，有限比例的粒子数处于 $\boldsymbol{p}=0$（凝聚）的状态.重新回到表达式(117.4)，应该（在化求和为积分以前）预先分出 $\boldsymbol{p}$ 或
$\boldsymbol{p}'$ 等于零的项，并考虑到这时在每一个 $\boldsymbol{p}=0$ 的量子态中的粒子数：
$n_{p=0}=\dfrac{N_{\varepsilon=0}}{g}$.此后，化求和为积分如正文所述，相应于(117.7)结果我们求得
\begin{equation*}
  \nu(r)=\frac{2n_{0}}{\bar{n}}I+\frac{g}{\bar{n}}I^{2},\qquad
  I=\int\mathrm{e}^{\mathrm{i}\boldsymbol{p}\cdot\boldsymbol{r}/\hbar}
  \bar{n}_{p}\frac{\mathrm{d}^{3}p}{(2\pi\hbar)^{3}},
\end{equation*}
式中 $n_{0}=\dfrac{N_{\varepsilon=0}}{g}$，并且 $\bar{n}_{p}$ 由 $\mu=0$
的玻色分布公式给出：
\begin{equation*}
  \bar{n}_{p}=\frac{1}{\mathrm{e}^{\varepsilon/T}-1}.
\end{equation*}
在距离 $r\gg\hbar/\sqrt{mT}$ 时，积分 $I=mT/(2\pi\hbar^{2}r)$（从上题中公式取在
$\mu=0$），因此，
\begin{equation*}
  \nu(r)=\frac{mTn_{0}}{\pi\bar{n}\hbar^{2}r}
  +\frac{gm^{2}T^{2}}{4\pi^{2}\bar{n}\hbar^{4}r^{2}};
\end{equation*}
如果 $T$ 不太接近 $T_{0}$（因此 $n_{0}$ 不是太小），第二项可略.在相反的情况下，距离
$r\ll\hbar/\sqrt{mT}$ 积分
\begin{equation*}
  I\simeq\int\bar{n}_{p}\frac{\mathrm{d}^{3}p}{(2\pi\hbar)^{3}}
  =\frac{\bar{n}-n_{0}}{g},
\end{equation*}
因此
\begin{equation*}
  \nu(r)\approx\nu(0)=\frac{\bar{n}^{2}-n_{0}^{2}}{g\bar{n}}.
\end{equation*}
值得注意，对于玻色气体，积分 $\int\nu\,\mathrm{d}V$ 在 $T<T_{0}$ 时发散，所以按(116.5)式计算就会导致粒子数涨落有无限值，这与§113中的有关评述一致.

\end{xiti}

\section{涨落的时间关联}

我们考虑表征处于热力学平衡的某个闭合系统或其一部分的任何一个物理量（在第一种情况下，有一些量如能量根据定义对闭合系统保持不变，所考虑的量不应该是这样的量）.这个量随时间有不大的变化，在其平均值附近涨落.我们仍然用 $x(t)$ 表示这个量的实际值与其平均值之间的差（以使 $\bar{x}=0$）.

在不同时刻的 $x(t)$ 值之间存在着某种关联；这意味着：$x$ 在某一时刻 $t$ 的值将影响它在另一时刻 $t'$ 具有各值的概率.正像在前几节中讨论空间关联时一样，我们用乘积的平均值
$\left\langle x(t)x(t')\right\rangle$ 来表征时间关联.在这里，仍旧像通常一样是在统计的意义上来理解求平均，即对于量 $x$ 在时刻 $t$ 和 $t'$ 的一切可能取值的概率求平均.在§1中已经指出，这样的统计平均等价于时间平均：在这里，在给定的差值 $t'-t$ 下对时间 $t$ 和 $t'$ 之一求平均.这种方式得到的量
\begin{equation}
  \varphi(t'-t)=\left\langle x(t)x(t')\right\rangle
\end{equation}
只依赖于差值 $t'-t$；所以，该定义也可以写成形式
\begin{equation}
  \varphi(t)=\left\langle x(0)x(t)\right\rangle.
\end{equation}
当时差无限增大时，关联显然消失，与此相应函数 $\varphi(t)$ 趋于零.我们还应该注意，由于定义(118.1)对于 $t$ 和 $t'$ 的交换明显对称，函数 $\varphi(t)$ 是偶的：
\begin{equation}
  \varphi(t)=\varphi(-t).
\end{equation}

我们把 $x(t)$ 当成时间的函数，同时也就暗指它自身按经典方式行事.然而，所述的定义也可以表达成适用于量子力学量的形式.为此，$x$ 这个量必须代之以量子力学的含时（海森伯）算符
$\hat{x}(t)$.属于不同时刻的算符 $\hat{x}(t)$ 和 $\hat{x}(t')$ 通常不可对易，因此现在关联函数必须定义为
\begin{equation}
  \varphi(t'-t)=\frac{1}{2}\left\langle
  \hat{x}(t)\hat{x}(t')+\hat{x}(t')\hat{x}(t)\right\rangle,
\end{equation}
此处的平均对精确的量子力学状态求\footnote{这里必须再提醒一下：根据统计物理基本原理，无论是用系统定态精确波函数以力学的方式，还是用吉布斯分布以统计的方式，求平均所得的结果不受影响. 不同之处只在于：在前一种情形下结果用物体的能量表示，而在后一种情形下表示成温度的函数.}.

假设 $x$ 是这样一个量，对它指定一个确定值（比它的平均涨落
$\langle x^{2}\rangle^{1/2}$ 大得多），我们就能描述一个不完全平衡的确定状态.换句话说，在给定 $x$ 值下建立不完全平衡所需的弛豫时间，假设比 $x$ 本身达到平衡值所需的时间小许多.这个条件对于很宽的一类有物理意义的量都满足.这类量的涨落称为是准定态的\footnote{也称为热力学的.}.本节下面将考虑这种类型的涨落，另外还假设 $x$ 是经典量\footnote{量子力学量准定态涨落的公式，最后可通过§124中所述的简单修正（参看(124.19)）从经典量的公式得到.}.

我们还将假设，在趋向完全平衡的过程中，在系统中所发生的对于平衡的任何偏离，都不需要另外引入新量来描述.换而言之，在每一时刻非平衡系统的状态完全取决于 $x$ 的值（更普遍的情况将在下一节考虑）.

设 $x$ 这个量在某个时刻 $t$ 具有比平均涨落大得多的值，即系统远离平衡.那么可以断言：在以后的时刻，系统将力图趋于平衡状态，与此相应 $x$ 将减小.同时，由于所作的假设，$x$ 的变化率在每个时刻完全取决于 $x$ 本身在该时刻的值：
\begin{equation*}
  \dot{x}=\dot{x}(x).
\end{equation*}
如果 $x$ 还是相当小，则可以把 $\dot{x}(x)$ 按 $x$ 的幂次展开，并只保留线性项：
\begin{equation}
  \frac{\mathrm{d}x}{\mathrm{d}t}=-\lambda x,
\end{equation}
式中 $\lambda$ 是正的常数；展开式中零级项不出现，因为在完全平衡态中（即 $x=0$ 下）速率
$\dfrac{\mathrm{d}x}{\mathrm{d}t}$ 应该为零.方程(118.5)是描述非平衡系统弛豫过程的线性化宏观“运动方程”（弛豫过程物理性质完全依赖于 $x$ 这个量的性质）.常数
$\dfrac{1}{\lambda}$ 决定建立完全平衡的弛豫时间的数量级.

回到平衡系统中的涨落，引入 $\xi_{x}(t)$ 记某量在它于 $t=0$ 时刻有给定值 $x$
的条件下在稍后时刻 $t>0$ 的平均值；这个平均值一般不为零.显然，关联函数 $\varphi(t)$
可用函数 $\xi_{x}(t)$ 写成
\begin{equation}
  \varphi(t)=\left\langle x\,\xi_{x}(t)\right\rangle,
\end{equation}
式中平均仅对初始时刻 $t=0$ 时 $x$ 的不同值的概率求.

对于 $\xi_{x}$ 比平均涨落大得多的值，从方程(118.5)另外有
\begin{equation}
  \frac{\mathrm{d}\xi_{x}(t)}{\mathrm{d}t}=-\lambda\xi_{x}(t),
  \qquad t>0.
\end{equation}
考虑到 $\xi_{x}(t)$ 这个量为平均值，应该认为这个方程对于任意小的 $\xi_{x}(t)$
也同样成立.把方程积分并记住根据定义有 $\xi_{x}(0)=x$，我们求得
\begin{equation*}
  \xi_{x}(t)=x\,\mathrm{e}^{-\lambda t},
\end{equation*}
最后，把它代入(118.6)，我们就得到确定时间关联函数的公式
\begin{equation*}
  \varphi(t)=\left\langle x^{2}\right\rangle\mathrm{e}^{-\lambda t}.
\end{equation*}
然而，该公式当前的形式仅适用于 $t>0$，因为在推导中（方程(118.7)）实质上假设了时刻 $t$ 在
$t=0$ 之后.另一方面，考虑到函数 $\varphi(t)$ 是偶函数，可以写出最后的公式：
\begin{equation}
  \varphi(t)=\left\langle x^{2}\right\rangle\mathrm{e}^{-\lambda|t|}
  =\frac{1}{\beta}\,\mathrm{e}^{-\lambda|t|}
\end{equation}
（$\langle x^{2}\rangle$ 由(110.5)给出），它适用于无论正负的所有 $t$
值.这个函数在 $t=0$ 时具有两个不同的导数.这种情况起因于我们考虑的时间间隔比建立非完全平衡（在给定的 $x$ 下的平衡）所需时间大得多.考虑不在“准定态”理论可能范围内的较短时间间隔，当然在 $t=0$ 时会有等式
$\dfrac{\mathrm{d}\varphi}{\mathrm{d}t}=0$，正如 $t$ 的导数连续的任何偶函数所应该的那样.

上述理论还可以表述成另一形式，它可能更有利.

如前所述，量 $x$（而非其平均值 $\xi_{x}$）的方程式 $\dot{x}=-\lambda x$，只是在
$x$ 的值比平均涨落大得多的情况下成立.在任意的 $x$ 值下，我们写
\begin{equation}
  \dot{x}=-\lambda x+y,
\end{equation}
这定义了新量 $y$.虽然 $y$ 这个量的振幅绝对值绝不随时间改变，但是如果 $x$
足够大（在上述的意义下），则 $y$ 相对较小，在方程(118.9)中可略.用这种方式引入方程(118.9)中的量 $y$（通常称为“随机力”），应该看成是 $x$ 这个量涨落的起源.这里的随机力关联函数
$\langle y(0)y(t)\rangle$ 应当适当限定，以使得可导出
$\left\langle x(0)x(t)\right\rangle$ 的正确结果(118.8).为此必须设定
\begin{equation}
  \left\langle y(0)y(t)\right\rangle
  =2\lambda\left\langle x^{2}\right\rangle\delta(t)
  =\frac{2\lambda}{\beta}\,\delta(t).
\end{equation}
很容易证实这一点，只要把方程(118.9)的解写成
\begin{equation*}
  x(t)=\mathrm{e}^{-\lambda t}
  \int_{-\infty}^{t}y(\tau)\,\mathrm{e}^{\lambda\tau}\,\mathrm{d}\tau,
\end{equation*}
并对乘积 $x(0)x(t)$ 写下二重积分然后求平均.

表达式(118.10)在 $t\neq0$ 时变为零，这一事实意味着：$y(t)$
在不同的时刻的值不相关.实际上，这个说法自然是近似的，并且只是表明：$y(t)$
的值只在与建立非完全平衡（在给定的 $x$ 下的平衡）所需时间同数量级的时间间隔内才是相关的，正如已经指出的，在上述理论中它是可略小量.
